\PassOptionsToPackage{quiet}{fontspec}
\documentclass{zhvt-classic-Miya}

\title{大至閣詩}

\begin{document}
	
	\maketitle{紹興~諸宗元}{三卷}[大~至~閣~詩]
	
	\insertgraphics[width=7cm,angle=90]{pic}
	\cleardoublepage
	
	\setcounter{page}{1}
	\tableofcontents
	
	\mainmatter


\chapter*{序跋}
余識貞長逾二十年癸丑甲寅間貞長官京師見輒談藝又時時相聚飲博越十年癸亥余居上海貞長參浙江軍幕其主將每招余至杭州暇則與貞長游湖上飲酒樓各出詩相視以爲笑樂又二年乙丑余在樞府邀貞長北來治官書晨夕相見顧簿書塡委而文酒之樂邈不可得居數月余謝病去貞長倉黃南歸以貧故復爲人掌書記體力渐渐衰退矣越六年辛未以病起樓詩寄余大連盡一册皆絕句讀其詩私心慨歎憂其不久於世是年余來上海復與相見輒和余妙高臺二詩未幾又病余走視諸寢因懷參餌貽之貞長目余曰環堵蕭然語次長喟余亟亂以他語壬申三月游華山歸貞長已前死三日夫以貞長文學粹美交友遍天下傭書老死而不獲一日之逸士之憂生失職至於此極然則詩人多窮之說其信然耶貞長治詩垂四十年不名一家而所詣與范肯堂爲近陳伯嚴鄭太夷俞恪士黃晦聞夏劍丞李拔可交口稱之余最喜其靜安寺追懷恕齋一詩以爲瀏亮沈痛而家國身世朋友之感胥寄於是蓋貞長嘗居湖廣總督幕府恕齋則總督瑞澂字也貞長才氣橫溢賦詩喜和韻和落葉詩疊韻至四五十不肯休朋輩無抗手者顧其過人處則在獨謠孤詠情與景融悠然意遠而不繫於更唱迭和之所爲也貞長旣死其友朱缽文爲之董理遺稿凡七巨冊余鈔得三百二十篇刊布於世稿本則歸諸缽文度必有好事過余而舉授之梓者然卽此以概其全亦足以盡吾貞長矣甲戌秋日長樂梁鴻志
\begin{fw}\end{fw}
余年十二隨先君至贛先君爲妙選師友故一時賢俊多得從之遊厥後聚散不常遭遇復各異曩者晨夕談藝之樂一逝不可復得日月易邁往昔盛年各隨塵浪以俱去三十年中哭桂伯華劉玉珩沈筱宜矣繼又哭陳師曾梅斐漪最近復哭諸貞長與劉未林文公達余幸稍識生死無常之理不至以此傷性然情懷則可知也又念諸君以特異之姿閱世數十寒暑所得者惟憂苦疾患此固人類所同然第求稍獲世間名利恭敬以竊自慰其非虛生浪死其事亦至復不易嗚呼人生之多艱文人之無命殆不信耶殆不信耶貞長長於余數歲少好爲詩屢相唱和驚其骨格騰健望之莫及中歲偶讀所作則益轉爲蒼渾駸駸與散原吷庵諸公並馳沒後覓所爲詩固裒然成集梁子衆異因獨任剞劂之役刊成此帙焉貞長平生知友甚黟今獨得此於衆異於死友可云不負抑貞長異時屢助節使幕府嶷然將有以自見固非欲以詩人終者頃歲胥疏無俚買地西湖紅桕山莊葺屋數椽期將隱居以老復遇焚如燬藏書萬數千卷都盡意氣乃頓衰謀食來海上所爲詩始率易及病以死幾無貲爲殮綜其身世頗類子美樊川與玉溪劍南然斂華就實能用於世復似過之至於詩之與古人絜長較短則未易以一二言盡也余往者曾屬友人吳眉孫刊師曾遺詩伯華以學佛棄其文字玉珩早死無遺稿筱宜斐漪所作散佚余比搜輯公達之作尙未成編未林遺稿則貧不能刊嗚呼當今之世貞長之猶得流傳此帙得不謂之厚幸歟抑文字與世諦牽纏至固貞長曾從伯華學佛今兹擺脫落塵鞅將舍其所執而企其所未至區區一詩集之刊布其未必爲貞長之意乎第朋輩之所得而助於貞長者則固止如是也噫民國二十二年十二月番禺葉恭綽
\begin{fw}\end{fw}
諸君貞長旣歿年餘其故舊取手書詩若干篇復輯錄朋游傳寫投贈之什彙爲若干卷長樂梁衆異遂擇其尤醇者推次年月編而刊行使余序之蓋君所爲詩黟矣然不自愛惜投積敗篋中草或不完己巳前存杭州宅者不戒於火君嘗謂我當憶存之今求其家不果有也君少才力橫肆好魏源龔自珍之學顏所居曰默定書堂中歲始更名大至閣所賦詩多隨興所至振筆書之未嘗刻意鍛鍊求勝於人詩成輒自喜方一牋置前未必盡愜吾意及彙百十篇誦之則才逸言雅固不在近頃諸能詩者後余交君逾三十年覽省其間酬唱之作若歷夢境悲憂患難嘗與共之有徘徊肝膈而不能去者況與君生同歲顧後死以序君之詩耶君善記書久客疆吏幕辛亥武昌兵變君遂去後過恕齋有詩最爲朋曹所傳誦生平志事亦略具於是又嘗記相者言當獲身後名紀之以詩翳古以來以詩名而命蹇似君者代有其人亦惟使吾輩執卷吟諷歎賞而已夫身後名果於君何濟哉新建夏敬觀序
\begin{fw}\end{fw}
病起樓詩一卷總七言絕句七十二首爲居士今歲病兩月餘時所製中間人事之乖迕病狀之嬗易皆著於篇什今幸而不死自視輒爲啞然但以詩而論則錢牧齋所謂許秀才體矣兒子章世祥世旣爲寫錄將以寄海內故人吾友朱翁炎午陸君丹林乃舉以付印海內之知居士者見此短册當信居士猶在天壤間也中華民國十九年四月二十五日大至居士記居士病之得生朱翁與張十一弟森玉皆有厚於居士用並記之
\begin{fw}\end{fw}
說明據民國二十五年梁鴻志所刊大至閣詩錄入原書一册不分卷收入爰居閣叢書卷端有散原手書諸貞壯先生遺詩時有重刊此書者以價昂未獲一覽王兄子衿不厭煩勞掃描見惠銘感無已顧青翎錄入
\chapter*{大至閣詩}


\begin{fw}乙未十月廿日得笙伯書\end{fw}

小雨愔愔晝閉關一庭冬翠靜如山殘年只合書中盡遠信今從江上還貧賤情親同性命風塵涕淚屬清閒危灘百尺袁州路爲爾還歌行路難

\begin{fw}感春\end{fw}

桃葉柳花今在眼荻芽蘆笋已登盤端居坐惜歲時晚狂語誰爲醒醉看殘雪壓檐晴亦雨高城近水晝生寒屈原幽隱韓孤憤未敢題詩強自寬

\begin{fw}晚望\end{fw}

高城橫暮氣煙樹晚冥冥江靜沈斜日天低耿大星疏枝時墮鳥暗草忽驚螢永念荷戈者哀笳不可聽

\begin{fw}私淑魏先生\jz{源}龔先生\jz{自珍}之學自題所居曰默定書堂爲此道意\jz{丙申年作}\end{fw}

龔魏文章數弟兄曾侯王叟舊相稱\jz{第一句曾惠敏贈王壬秋句}抗心志在從人後邀客書成榜我楹百歲未過多變駭一言可蔽本縱橫人間紙墨難求得\jz{謂二先生著述}誰與狂生訴不平

\begin{fw}檢東方朔傳上書年二十二與余今同作自詠一首\end{fw}

亦讀人間有用書中原㼱\jz{㑺耎二音}滑氣難除還同臣朔年如許自責何能得自書

\begin{fw}戊戌秋中寄東虛\end{fw}

聞汝全家去撫州未將消息告交遊人言張儉同亡命我識梁鴻豈避仇有舌尙存言變法此心未死莫悲秋種瓜自是君家事莫向青門問故侯

\begin{fw}湖口舟中望石鐘山\end{fw}

揮手匡廬蒼翠間渡江先看石鐘山荒荒雲木明吾眼歷歷亭臺擁此鬟往歲記曾笳鼓競少年我媿鬢絲斑伏波橫海今銷歇擊檝他時幾往還\jz{謂彭剛直}

\begin{fw}感紀\jz{辛丑年作}\end{fw}

柳絲池館過清明庵牓好題心太平地僻近多疑遠信天愚寒不解朝晴草間求活留花種雲外遺音有雁聲不信南東金粉地供人揮涕說神京

\begin{fw}竊憤\end{fw}

亦許天光一畝開星球九萬里雄哉不知有漢桃花老旁若無人燕子來竊憤無端供涕淚遠情何事苦喧豗松陰幽室求文稾猛士終題廿一囘

\begin{fw}暑熱\end{fw}

不眠先鳥起破曉有蟬聲樂歲今憂旱驕雲尙喜晴窓明蟲闇聚廊靜鼠橫行睠此清涼境昏昏惜衆生

\begin{fw}伯華師自日本來書云近與吾友通州范彥殊彥矧兄弟相唱和旣以書報賦寄長句\end{fw}

觥觥德化桂夫子更念通州兩范生日共哦詩對東海夢憐羈客在秋城脫身幸自兵間出盡室今爲浦上行欲補春秋告三世直從據亂到昇平

\begin{fw}同友人過味蒓園\end{fw}

夾道林光受雨多高雲不動亘天河明明釵朵當鬟見一一車燈啣尾過難忘臨筵歌定子莫妨歸笑語黎渦虛廊燭暗還吟嘯其奈清宵百感何

\begin{fw}夜遊約同去病\end{fw}

酒餘奇氣不能收林薄驅車作野遊片月早升光在樹萬家無睡夢成秋樓臺近海宜忘暑士女圍鐙不解愁難遣近來懷抱惡悽然越謾復吳謳

\begin{fw}移居吳下感賦\end{fw}

已忘門外卽天涯僦舍初安靜不譁太息賃舂今失婦豈徒家具少于車越風漸近當懷土吳俗能知可戒奢兒髻嫗裙我書卷村夫子欲畫移家\jz{村夫子移居圖宋人有此本}

\begin{fw}吷庵以五言一篇慰我婦姚怛化之戚引其昔感造詞甚悽且約近嘗爲七言歌詩以洗近體嵬詭卑俳之習因依韻變體畣之楞嚴經所謂譬如飲水冷㬉在心可共喻也\end{fw}

微生分合戰憂患至人長欲齊彭殤兒時發願今已壯朱顏不駐翻成蒼去年哭父今哭婦直如決疣生奇創得從地下見翁媼勝我百感摧肝腸夏生顧我長太息爲言舊厄曾親嘗虛帷鐙閣見魂魄是夢非夢疑哀荒\jz{君語我其前夫人逝時其狀如此}詩來慰我雜悲哽坐忘斜日明池廊嗟哉貧賤誰相依對泣猶念牛衣旁

\begin{fw}九月十一夕乘月自閶門歸過王廢基感賦一篇\end{fw}

寒月潑水成空明垂柳結幄馳道橫言從西城趣東郭長衢短陌多經行兹來胡爲獨感喟陳蹟了了思隆平\jz{元時蘇州爲平江路士誠初據時改爲隆平}奇渥溫亡賴梟桀崛哉張氏爲雄爭吳民至今尙尸祝七月晦日香花迎大起園林盛館第如何委墮無遺甍姑蘇臺荒糜鹿走徒使千載存其名改元定歷正官制一時海上連諸城功雖不成志足喜白駒場人張士誠西風獵獵黃菜葉坐棄天下猶可驚廣場三里今何用留與後人供校兵

\begin{fw}伯嚴丈招飲淮榭明日賦謝\end{fw}

枯淮得雨已通潮避舸移鐙坐此宵倦客因花多有感溼雲籠月不能高俊譚信可驚流俗短簡還應謝盛招捉扇登車了無怍我猶難解李生嘲

\begin{fw}飲歸\end{fw}

街泥妨轂雨懸廊爛熳還憎歸路長但覺衫裳携酒氣閒從牕壁數鐙光君將南去憐羈旅歲有秋風感莽蒼神駿何能賞鷹馬賸持缾盋禮支郎

\begin{fw}午日同劍丞遊留園晚飲市樓缶老約也卽事有述\end{fw}

酒場吟地已非春重午來尋九陌塵園外遊船時見伎竹邊籠鶴尙親人佳晨翻笑閒成病\jz{時病足}晚醉能敎句有神呵壁更忘天可問何因強起屈平身

\begin{fw}漚尹先生雨中招飲聽楓園席間舉東坡事以爲笑樂歸成此篇亦申其例\end{fw}

好客經旬再舉觴照廊花樹粲成行吳中高會今難兩酒半悲歌意更長往事古今成俊語虛堂鐙火有雄光東坡樂府傳新刻我向黃州到武昌

\begin{fw}爲婦姚營葬於南昌鄧家步先塋之次感悼成一篇\end{fw}

臨穴猶能一撫棺離離生死語原難但期後日同歸此忍說人間百不懽合眼庭堂記歌笑傷心先隴伴荒寒伯倫酒德今誰頌落日高原荷鍤看

\begin{fw}庚戌三十六歲生日賦示親知\end{fw}

昔我年十歲讀書方閉扃文選能上口座客稱甯馨甲午越二十自壽齊歲星晨筵張庭堂晚醉傾盃甁結交多少年豪語輒可聽慷慨言憂時糠粃視六經迨年三十時喪母嗟零丁寒舟檥吳城獨登湖上亭眼中五老峯對我同冥冥日月忽焉逝卅六成今齡顏枯難再澤髮縞難再青託身慕冰蘖息念忘風霆豈無稻粱謀但乞文字靈黑頭作卿相黃金聚娉婷此心久寂滅無待書其銘醵飲勞羣公會食風中萍何以厲吾志何以葆吾形還當乞嘉言不寐先求醒

\begin{fw}飲拔可家時予將返吳矣\end{fw}

對岸相看江作峽到門今藉橇乘泥羈懷寒與千峯遠客座忻同一笑齊到眼新髭雜衰鬢稱心白酒與黃齏子留我去銷殘臘短詠長謠約共題

\begin{fw}拔可書來屬過江甯一視其家往訪不得門巷賦示拔可\end{fw}

來訪谿堂戒夙期出門惘惘路成岐迴車曲巷終相失高柳叢花我所知此諾竟敎呼愧負故人應不爲瑕疵市橋小立風吹袂衝雨歸來賸有詩

\begin{fw}與海藏相遇京師有詩見贈賦答\end{fw}

牽率無端入國門庚申庚子蹟猶存君思百日維新夢我有千秋舊涕痕舉世殆沈三里霧說醫難洞一方垣鍵關人海終歸去自惜狂言舌可捫

\begin{fw}辛亥元日試筆\end{fw}

吳中吳中三年居吉利吉利元日書有妹得壻嫗抱子它無所樂無所須閉門不出坐聽雨稚女新衣作嬌嫵人生富貴亦等閒情親祗樂長相聚

\begin{fw}辛亥五月武昌苦雨之什\end{fw}

五月一日雨人謂將熟梅昏昏雲接水宛宛壁生苔占竟經旬驗聲何倒海來誰爲桑土計床漏未容哀

\begin{fw}芝罘舟中贈澄一\end{fw}

船廊握手語匆匆魯酒難令客頰紅一别秋随人在遠此行天使海無風浪波坐付沄沄歎蹤蹟翻憐處處同歸後蒼茫有何意盡拈書卷問牆東

\begin{fw}初渡海至芝罘賦示拔可\end{fw}

平生未與海相狎此去能看日出東今與老坡分一席登州樓觀得毋同

\begin{fw}夜坐示拔可度公\end{fw}

滅燭窗虛見海光瞢騰但以睡爲鄕風中巨浸知山立\jz{錄者按知疑當作如}秋後清宵比晝長三客復成今日聚百端難諱少年狂輕寒疎雨簾垂地同有南人念北方\jz{時眷屬均在南}

\begin{fw}簡介庵\end{fw}

同官如所親好儒或多緩才劣身則勞意廣力不滿賢者不忍棄往往縫我短囘頭已隔世袖手入孤館看花曉覺寒及物春不暖沈冥誰者豪千載魂自斷

\begin{fw}雪晴後賦寄吷庵\end{fw}

雪晴頗喜市無塵雲暖還疑歲已春列榻罽茵宜睡美浮杯花乳試茶新朋交百輩留江介蹤跡何人到海濱已勝東坡在儋耳和陶詩卷直隨身

\begin{fw}同竹君拔可求所謂小蓬萊者至寺門中駐兵戒不入小立松下遂歸其地可望海且攬全市形勢殊勝絕竹君云不可無詩余與拔可遂賦紀之\end{fw}

山如奔馬不能羈到海還留一境奇飽聽松聲竟歸去未知蕭寺建何時孱兵尙戀僧廚粥過客誰題石碣詩欲問浮雲望西北近聞胡騎正東窺

\begin{fw}癸丑春題太夷海藏樓圖\end{fw}

收拾堂堂意斯樓築海濱插天環堵宅絕代閉門人種樹春多異無山境自新讀經資暇日惜我不爲隣

秋晚賢良寺相從慷慨言救亡終一策奮舌有羣喧惟我能相喻嗟今不足論四年騰踏過攬卷夢飛翻

\begin{fw}夜過海藏樓歸紀所語簡太夷幷示拔可\end{fw}

海雲歛月涼生芒栝竹影地森成行馬蹄蹴踏若翻水但惜損此寒瓊光打門夜半撼鄰睡主人延客先下堂兼旬再見已足喜况能坐對秋宵長主人論詩得眞理近稱米氏推襄陽上規韓白許永叔出以簡淡非尋常老顚落筆不局曲意境往往齊蘇黃前聞沈侯語夏五文字不可輕其鄕此言當爲永叔發我意若合衡與量嘉祐元祐在何世縱有作者人謂狂明星耿耿露作霜我歸所止神旁皇人間俗論成嗤點段拂豈果師元章

\begin{fw}同吷庵來鎭江居金山下曹氏園用東坡兩賦金山涼字韻索吷庵伯蓀栗長和\end{fw}

犖确車輿度夕涼到門草樹雨餘香去年離亂今初定往日經過客笑忙\jz{十年來過金山未一登眺}秋與閒時逃寂寞天圍山氣作圓方塔燈照水鐘鳴夜忍逐吾懷墮混茫

\begin{fw}園池上有廢亭同伯蓀坐話久之次前韻\end{fw}

園柳青疏晝亦涼渚蓮紅退葉能香是間濠濮閒成趣我輩蛩驅老尙忙已負買田爲下策頗思種樹乞奇方謀人家國原餘事難得山川闢莽茫

\begin{fw}畣吷庵戲調再次前韻\end{fw}

環溪廊屋早生涼竹粉苔衣不斷香小住當爲開口笑故人爭累和詩忙宵中酣寐吾何適佳處茅庵古可方獨惜高眠失清景江聲推月破昏茫

\begin{fw}同伯嚴丈味蒓園茗坐\end{fw}

車馬來稀我到頻日斜廊靜見窓塵眼中過去莊嚴刦座上無多磊落人抵几惟聞談戰事煑茶隨分著閒民百錢尙可黃壚醉亂世爲儒未足貧

\begin{fw}吷庵以與趙堯生楊昀谷諸公叠和詩見示趙蜀中詩人去年秋曾於朱漚叟所見之吷庵和詩多屬漚叟叟今居蘇州予亦不能無憶也遂依韻寄懷幷示吷庵\end{fw}

聽楓園讌秋雨過座客論才尊俎左今年能對去年人憶昔相逢足相賀風蟲日鳥青春深照水緋桃新萼破縱無嘉會相招邀賴有清篇騰咳唾獨思漚叟居吳中烽鏑不驚詞作課偶來海上輒急歸幽憤何殊老雪个近狀如我聞相同家人病減除藥銼周廊花樹張春晴惜未涪心寮裏坐此間日見惟夏五天何不仁多坎坷夏喪厥妃我哭妹哽痛哀呻相繼作人言愁苦踰邱山獅象難移况蟻馱我曹性豈甘愁苦夢飽但期醒覺餓詩成不敢貽朋遊手寫惟趣朱夏和

\begin{fw}園游賦示孝若\end{fw}

林陰草色漲晴流不雨來登池上樓十日塵埃閉亭館一年花事有春秋故人何暇來相慰壯歲多憂強自休橫笛按歌吾久倦南歸今始作園游

\begin{fw}過吷庵新居\end{fw}

跨河隱高橋導徑歷深樹君家近移此我已四來去主閒客自適僮熟語不誤時時作歌嘯往往達曛暮前遊陟園館車馬遊馳騖昨出踏衢巷巾襪滿塵土信知腰腳健足脫形骸苦悄焉倦言歸一夢謝百慮注水鳴窗廊到耳夜疑雨知君不能眠詩成漫相許

\begin{fw}和拔可韻示度公\end{fw}

唫苦詩寒謝俗紛未應好事獨輸君客閒自喻寥天鶴樓迥惟看度海雲能與狂歌送昏曉非關禪意斷知聞寒宵跂腳吾何語細數歸鴻正作羣

\begin{fw}散原以庚辛年間詩卷見示感上一律\end{fw}

一卷冬春上塚詩近嗟節日廢歸期憂天淚盡翁垂老避地人多事可知宵夢長敎戀松檜餘生未分狎蛟螭朱姑橋近梅家巷\jz{此二地在南昌城南先塋在焉}麥飯無由補我悲

\begin{fw}吷庵示移居詩賦畣\end{fw}

昔君賦移居惟我曾投詩一年三徙宅君我胡同之歲晚得幽築位此詩人宜廣園有十畝下復臨清漪春風扇萬綠鸝燕來相窺閨中婉孌婦春及親迎期君其學高柔畝川甘所師陳哀鍥不舍吟諷何酸淒不爲新人笑乃爲舊人悲遂視棟宇間一一摧心脾我欲廣君意哀樂同嬰兒墮地一聲哭憂來知無時舉家今繫君甯能徇其私况兹天地閉海內紛流離一室事掃除我久爲世嗤幸君強自寬毋忘語者誰

\begin{fw}雨中賦答拔可依韻\end{fw}

吾生在野早聞道下士相逢先絕倒與君握手十載強誰驅白日迴蒼昊操竽終不變齊瑟奮弩何能穿魯縞人材令僕正紛紛賸可閉門同却埽北居五月今南來坐負花時被花惱家無半畝終安歸心憂一雨遂成潦漪園昨過花如塵賴汝新篇供饋犒添丁待抱爲我祝患難相期得自保卽言任土治貧瘠但數驚雷起枯槁平生自許果如何莫遣紅妝嗟我老

\begin{fw}夜從靜安寺道歸過恕齋故居聞恕齋昨日葬西山矣感紀一篇\end{fw}

故人築宅臨道邊我亦居此曾三年明明七載同所止吳風楚霧心茫然西山昨日成新阡一棺戢身碑有穿園中草木皆故物衞公戒子悲平泉人間一事猶掛眼當日屋成君病眠憂時憤俗復強出篤守故轍無他遷翦除豪滑奮彈戮時論百喙稱其賢惟剛及介君可信宮府不通繒與箋國門再入騰謗譽言有不用憂悁悁鳴笳建纛送還鎭我時竊歎君難全牙兵日橫柄濳奪崇朝禍發張空拳古來覆敗類如此爲之者人成者天江聲東流日夜急君何負此清冷淵黃州鼓角不可聽别語雖祕人能傳戶庭重過見寒月榆柳已老無鳴蟬迴車腹痛在何日我今躑躅來門前

\begin{fw}贈畹華\end{fw}

三年不見汝如何移笛桓伊白髮多絕憶丹邱珍重語南人爲曲北人歌

格勢誰從樂府知繁絃急管囀喉遲砑光帽畔山香舞曾對風軒落帽時

缶翁近爲寫疏梅貌取清閑亦費才綴玉更題軒榜字誰呼白石老仙來

\begin{fw}題畹華小影卽以贈之\end{fw}

目論紛紛說賈馮但宜箸錄付明僮如何絳樹雙聲者今在黃塵十丈中孝潔頗聞關至性品題端可謝羣公年來已嬾評箏笛淵默雷音許強同

\begin{fw}沈濤園丈招同癭公秋岳嗇翁酒集鳳卿畹華亦在座賦謝一篇\end{fw}

酒盡休言客已慵月高始覺暑無風趁閒料量娛今夕譚往蒼涼得二翁捉扇看求書跁跒聞歌醉倚喚玲瓏五雲四芷何銷歇賴有王梅在座中

\begin{fw}拔可補寫二十年所爲詩病中起讀不覺竟卷賦示一篇\end{fw}

一病遂敎弛百骸君詩娛我晚來佳舊唫辛苦爲追寫此事從容有好懷學古自能成獨到嗜奇不信賸吾儕西江宗派今無幾坐被人稱近簡齋

\begin{fw}畣吷庵幷示拔可\end{fw}

平生寡結納所交非常人曰有夏與李爲我平生親三年海上居遠俗羞爲隣高蹈匪本懷探手謝雜賓近役從檢校惟李同宵辰遂惜夏生遠汗漫忘青春團扇不相窺君言别浹旬車輪可日數自辨甘勞薪羣暌苦相憶詩來無辭頻和詩久不屬太息喪吾眞

\begin{fw}重九賦寄寐叟缶翁\end{fw}

嵯峨雲物伴沈吟豈有人間辟惡心佳日先知少風雨故鄕端可罷登臨授衣霜後棉爭市脫帽花前菊滿襟漫更裁牋告陽九此愁無古亦無今

\begin{fw}晦聞云有詩極寫近日荒嬉之狀余未得見也雨坐遣懣亦成長句\end{fw}

相對尊前浩蕩思不輸花底醉眠時清歌豈暇求車子倦眼何堪問偃師已信琴聲懸羯鼓偏從巾角論彈棋中山千日沈沈在戒飲無端近酒巵

\begin{fw}京口道中有寄\end{fw}

重有江南山照眼此行不樂過金焦杏花已信年當嫁桂樹知無隱可招晼晚風鐙曾照夢迢遙江水可通潮聞歌早悔輕言别誰解遲歸慰寂寥

\begin{fw}送人還京口\end{fw}

喜汝能歸更悄然鳳城南陌夢如煙心知一諾平生負事許千秋後世傳别日可思無強憶衆人相遇復何賢獨憐茶癖差同我好試中泠第一泉

\begin{fw}吷庵南還於天津濟南道中均有詩見寄君得歸而詩意多悵失若予者更何言哉依韻寄之並示拔可\end{fw}

浩蕩君今禮嶽歸心知憎與故人違却憐竟日搴帷坐定有低雲挾谷飛客眼醒時忘昨夢河聲盡處見斜暉冬春再别顏同槁豈獨衣裘減帶圍

卻曲迷陽君去住但能縱飲傚平原嗟予亦念盤中曲避世難尋瀼上村久信求田爲上策卽言乞食向何門約堂相見當相憶倒海翻河注酒尊

\begin{fw}紀别\end{fw}

鏡檻何緣貯淚光背冬涉夏去堂堂欲忘哀樂談何易苦計宵晨别更長平子未遺青玉案盧家誰築鬱金堂鐙前萬語終何意可作人間卻病方

\begin{fw}兒邁生賦示家人\end{fw}

阿翁尙戀爲兒樂舉室相呼喜汝生待付詩書作年穀不期愚魯換公卿得男雖晚吾何憾繩祖當如世所稱北發明朝更南憶高樓鐙影有嗁聲

\begin{fw}甲寅歲除前一日集法源寺爲陳後山先生設齋供先生歿於宋建中靖國元年卽是日也旣歸賦詩紀之\end{fw}

法源古蘭若盛集惟花時行廚沸酒炙下及屠沽兒今來歲殘臘竹凍無新枝邈焉千載懷酹酒將爲誰欲呼古今人一一皆詔之吾曹何好事獨爲無已悲生死本乘化時日無可疑豈如緜山火寒食哀之推座中有何言羣客稱其詩精靈亘天地餔糟與歠醨如身觸暑寒造化不敢私酒罷各散去吾悲復誰知

\begin{fw}元日出遊\end{fw}

城雲得日始能溫苑樹環煙靜不喧行散漫爲新歲樂邀頭猶數故風存將看燭影開馳道眞聽車聲過闕門鐘鼓爰居了何意經過鹿柴我無言

\begin{fw}缶翁喜余至海上以詩見贈依韻\end{fw}

臥憎髀病方春始來問樓居復臘初視昔交情若醇醴感翁文字爲爬梳皈心白業猶耽酒揮手黃金儻荷鉏尙可姓名相爾汝勝它李志與曹蜍

\begin{fw}缶老齋中懸自寫菊花幛子二奇橫之氣直可爲花寫照欲取而誼有不可欲置而吾又不能忘懷因爲此詩以投缶老如能舉一爲贈則吾詩不虛作矣\end{fw}

荒齋苦無菊秋色去大半凌晨登君庭菊也何燦爛異哉霜下傑不信出君腕何假甁與盎何待種與灌突兀千萬華照眼列枝幹落筆無醜姿觀者駴且汗信難凡卉同足爲吾見罕西風在籬落擷之不敢玩有如拜家姬雖美豈容讚狂思勞點染倦恐累閒散知難千絹酬敢託一詩換眷言癖斯堂\jz{缶老堂名}無以答客難

\begin{fw}雨晴出遊先投缶丈\end{fw}

窓日斜明風倒吹雨晴已信出遊宜曠然爽氣成秋矣來過閒坊問癖斯巨幅墨花容我乞古牋獵碣許誰知想公縑素如山積換取諸生七字詩

\begin{fw}和畣晦聞\end{fw}

著我疏鐙密幄邊了無情話亦翛然經行但遣中年感聊浪爭如六博賢燕子春鐙誰可喻馬蹏秋水不名篇紅裙豈復關文字一任他時好事傳

\begin{fw}雜書貽拔可\end{fw}

嚴車一日通夢已握君手纚纚千萬語張眼失唯否卽言在睡中了不異醒後世變未足奇可語在吾口今來暫相見别去願相守莫謂鈎有鬚或憶柳生肘樓居俯海甸愼莫酌我酒花枝大合圍黑頭看老醜

\begin{fw}南歸\end{fw}

夢隔高城送别人行行寒霧擁奔輪風塵直悔勞行役花藥猶能及早春家寄江南非萬里書言兒病已經旬歸來置膝將何語我更今年白髮新

\begin{fw}車過徐州感賦\end{fw}

扼險何人道已通川原猶說此邦雄疏星澹月隨車後殘角危城在眼中鐵騎防淮時已異羽衣吹笛惜難同一流宛宛趨山去豈是當年百步洪

\begin{fw}四月三日哀邁\end{fw}

汝生值今朝死先十五日生年未及周委化復何恤有母今哭兒有姑今哭姪爲言墮地事直至汝嬰疾我雖強不悲淚欲奪眶出幾忘殯在野忽念抱置膝我初聞汝病得書恍有失南還促宵征入門不敢詰見汝能笑嗁若脫械與桎無何病中變不謂救無術沈沈汝病時寒雨挾颶䫻晴光晚在牖汝竟一生畢我今寫詩處卽汝死之室魂小將何依夢囈雜啾喞

\begin{fw}送歸南者\end{fw}

君復南歸我獨留眼中雲物正清秋江喧風雨成荒漲民厄流離念本州且置買田陽羨約能無橫涕監門憂迂儒箝口論封禪賴可長飢對白鷗

\begin{fw}送栗長入蜀\end{fw}

頗聞灔澦已無灘君去休歌蜀道難豈似陸游耽晚宦莫敎嚴武笑儒冠中原多難憂戈甲十日能留戒雨寒玉壘浮雲有今古驛書題寄遠人看

\begin{fw}視殯\end{fw}

荒殯遠尋來視汝斜陽已落尙親人闔棺遂鐍三年淚了刼終留一聚塵自詭姓名原可嘆卽言生死復何因分明前病成今讖不到江亭亦愴神

\begin{fw}乙卯中秋\end{fw}

纖雲埽盡露華明此意何能語太清一夜思量在鄕國四年南北問陰晴廣筵燭燼香還爇夾巷歌喧月漸生此境兒時不堪紀未應扶夢坐深更

\begin{fw}徐州\end{fw}

北望寒雲曉不收南歸今始過徐州日光將出雲奔馬風力初溫渚泛鷗地畫中原風尙僿民居山磧氣難柔半生乘障思爲法世論悠悠孰可謀

\begin{fw}雨中夜發上海曉晴達金陵復渡江趨浦口舟中感紀\end{fw}

一夜奔車越六城昨行衝雨曉還晴亂流單舸浮江去障日千峰出霧明春遠初看榆柳大民閒猶說鼓鼙驚東南物力將何策悔不歸營谷口耕

\begin{fw}拔可旣卜居意有不樂賦此廣之\end{fw}

君昔卜宅金陵居自謂鍾山落吾手一朝捨去不卻顧久閉春牕與秋牖三年海上君獨閒十畝林陰今始守騰書前日告京師我爲申眉開笑口南來扣戶初不知突兀高樓隱榆柳開軒衆綠相周遮直疑鄰亦爲君有平生自許策無失此事何關狂以狃當前快意豈偶然草創中興聽評剖

\begin{fw}丙辰三月移家返杭而心白卜居海上心白索詩寫五言一篇寄之\end{fw}

僦屋安妻孥君其長子孫閱世重憂患君其弢語言悠悠吳楚郊往事安足論吾閒且歸居國敝將圖存奮身犯百難自歎無弟昆不如君有弟奉母資晨昏樓軒語宵雪海舶觀朝暾昔遊在夢寐苦語相寒溫吾居近湖山君可來扣門

\begin{fw}海藏樓觀櫻花投太夷\end{fw}

每歲花時輒詣君林光樓影漲春雲已看種樹三年大頗類爲香百合熏醉後晴陰難料量眼中朱碧強區分能分海氣來籬落歌哭人間可不聞

\begin{fw}泊下關傍夕大風雨雹\end{fw}

晚涼飛雹隔江看船尾風旗颭怒湍已信官蝗當漸戢但憂后土不能乾捲帷向夕天無暑破柱驚雷夢可安絕似五年前過此鍾山沈霧水平欄

\begin{fw}二月九日靈峯寺看梅遂過烟霞洞則梅已半落矣是日同游爲吷庵旭初\end{fw}

我昔知西湖以梅甲天下故者摧爲薪新者僅盈把荒區能久留天意不容假兹游同慕此有客曰汪夏北山踏芝塢靈峯扣蘭若飄香三百株柯萼森厏厊積縞俯泉素叢丹列巖赭吾行目爲眩陟亭倦腰髁\jz{寺後有亭曰來鶴是日未登}還徑趨烟霞餘春已傾瀉豈因來游人日夕相撏撦頗聞寺僧言利不及楸檟歲歲惟花時過此憎醜䰩踞廊石爲溫意不樂他捨何當遣東風一夜換嬌奼惟佛高出世如士遜在野此花將毋同忍使媚觴斝重尋期隔年書此告遊者

\begin{fw}同散原觚盦仁先游雲栖還賦一詩\end{fw}

入寺喜聞鐘磬聲二翁一客共郊行竹光飛雨能相待山勢趨江到此平林壑插天容衆綠澗泉緣徑識新晴此中獨背西湖勝壇席人間未許爭

\begin{fw}大至閣\end{fw}

隣畦判南北雞犬無甯聲所憎妾婦愚但爲箕帚爭

睡聲能撼眠近在屋前後中宵兩眼開不如繞床走

千金懸國門不得易一字乃聞說士雄輦金還自市

日不沈於西月不升於東升沈果何物遂旌宵晝功

\begin{fw}檢張幼樵澗于詩集\end{fw}

敗諱枋頭事豈然卻於詩卷見清言謫居強喜同和仲甥館幾聞累茂元往獄猶留諸將恥清流今有幾人存陳濤一役哀房琯士奮功名忍並論

\begin{fw}漫興\end{fw}

果能赤手縛於菟始許人間作丈夫我已安排他日事中年爲俠晚爲儒

\begin{fw}車行晚興\end{fw}

低煙落日黕平蕪道左新看柳萬株入夜車聲仍北去踰淮山勢尙南趨嗟予久墮江湖興徇俗難爲盎釜謨一飽終當念藜藿此時呼粥出行廚

\begin{fw}畣劉三見懷\end{fw}

落落吾將審所歸悠悠君與俗相違鍵關高臥非沈隱置驛來詩亦累欷小飲尙思秋禊樂狂歌不畏世人譏阿蠻醉倒當年語好事流傳涕欲揮

\begin{fw}舟出吳淞\end{fw}

一年三命檝行役亦尋常帆影奪空白江流受日黃澄清非可俟詼詭意難忘列嶼沈兵氣吾憂接混茫

\begin{fw}同漚尹病山觚广蒼虬遊雲棲賦寄伯嚴丈江甯去年秋曾與丈同遊也\end{fw}

環亭竹木春無改拍岸江濤我再來能與一山爲主客曾攜百感此徘徊社中白髮懷游躅宅外青谿照酒杯秋已隔年寒在眼寺門應憶翠成堆

\begin{fw}秋日雜興\end{fw}

昨雨今晴不可知昨晴今雨復何思閉門晴雨無關我但記移花抱甕時

\begin{fw}歲臘臥疴拔可吷庵存問稠叠心白復來杭省疾賦詩爲謝\end{fw}

病中不覺過除夕海內相憐有故人急欲報書知未死何如裹飯恥言貧閒敎宵寐能連曉靜喜冬寒已換春梅正作花吾彊起卻勝悵臥問漳濱

\begin{fw}輓沈濤老\end{fw}

強醉淵明漉酒巾疲轅往轍記轔轔棄官尙積回天志有壻曾爲殉國人\jz{謂林晚翠}舄奕功名喬木在崢嶸文字故楹新公原無憾時同惜雪涕諸孤語苦辛

雄篇晚勝復工書能事恢恢德業餘何意西風聞弔鶴應敎南徼激湍魚低回病榻猶談戲配食先祠合掃除生卒紀年皆戊午龍蛇歲厄竟何如

\begin{fw}觚庵旣歿之五月昨忽入夢追輓二詩\end{fw}

壯推健者晚蕭然絕類龍川及樊\jz{借作仄}川餘事能詩無近習積憂遂病損宵眠曾官隴外爲時重賸有荒亭可畫傳合眼廿年歌笑地忍敎留語閱桑田

君隱湖堂我亦歸但驚今瘦遜前肥巷居自僻勞經過妹服初持感累欷\jz{余初居慶和衖君時見過去年十一月九日余有張氏季妹之喪君猶走慰距其殁僅十二日也}槃食漸耽霜後菜寺游回扣雨中扉\jz{此二語皆紀君歿前二日事}如何入夢忘生死爲道年來不解饑

\begin{fw}同劒丞步湖上公園乘月歸賦成索和\end{fw}

四年歸對山與湖一日不捨朝與晡朅來寒霞未西沒仰見月出東北隅人言開歲過四日何用舊歷追望舒今月古月本無異賴此寥寂能相娛樓扉多鐍冷墐戶鐙竿相互明綴珠山光隔湖不入夜宛宛盡向城中趨夏侯及我乃耽此但惜園卉多枯株圓靈燭霄未知紀清游狂賞無時無往者躑躅今喁于會合天畀非須臾投篇世必駭光怪隔巷日可同歌呼徐亭蘇圃忽相憶攜家後亦曾居吳當時躍馬兩年少君今白髮吾白鬚

\begin{fw}孤山晚興\end{fw}

寒同閒客坐斜曛打槳清波靜可聞出嶺雲偏留一抹界湖隄自作三分花光娛我遲當見塔勢如人了不羣欲寫屋山忘歲月爲招顏寓及蕭耘

\begin{fw}鎮江至高郵日再易舟\end{fw}

偃蹇家居久經行夏雨寒車中共誰語江北少山看分灌淮方涸趣程道更難輕舟日再易直似上風灘

\begin{fw}寶應道中賦寄嗇翁\end{fw}

隄高橫巨浸潦落苦驕晴匡坐依帆影來船問馹程過江非寂寞行水愧經營廿載治淮略何能慰友生

\begin{fw}晨起聞雪甚盛\end{fw}

屋後風篁折凍聲我知急雪到天明林昏湖淨勞宵憶罏煖裘溫記昨晴多睡恐來僵臥誚獨遊期與苦寒爭吾文今作田夫禱螟螣今年不再生

\begin{fw}正月二十日同愔仲仁先集煙霞洞梅方盛開遂用東坡是日女王城詩韻約同賦\end{fw}

迎客孱僧不閉門煙霞石屋碧連村來從荒壑尋春事下視江流有漲痕柯萼任攀香未減巾裘待換午能溫去年亭柱題名在了了同遊役夢魂\jz{去年正月二十一日同拔可秋岳衆異來游}

\begin{fw}湖上春行各繫以詩\end{fw}

四百五十樹東風寺閣花光薄日烘何似我園春一樹高枝可敵衆株紅\jz{靈峯寺看梅}

拊檻呼魚撥剌鳴堂堂策策本天成誰敎豢養移根性投餌甘爲大小爭\jz{玉泉寺看池魚}

\begin{fw}得癭公京師書卻寄\end{fw}

城中有山朝暮見紫翠晴昏看百變出門繞屋俯谿流已與澄湖鬥飛練南居勝北此語久家食累人吾計倦京師一月書再來爲說花時共游醼故人磊落惟羅癭大隱王城謝時彥平生所交徧海內識我初知我狂狷索詩寫寄還悵失在口能吟不如面朅來江上多北風過去人間怨漂霰大隄士女唱銅鞮誰問襄陽耆舊傳山川信美况鄕國捨此他求復何衒撑腸文字閉門居我亦勸君勤筆硯

\begin{fw}晚行湖上記與友人話\end{fw}

照水鐙竿遠受風城隅非月亦溟濛得詩仍在湖光裏著我何能野艇中一雨春衣寒又換萬喧歌市夜方空頗聞今歲新絲賤誰謝桑蠶作繭功

\begin{fw}湖上雜書\end{fw}

疏楓炫晝光隨日淺渚安流氣作霜坐笑龍川喜風雨彘肩斗酒渡錢塘

先生顏貌漸如翁湖氣隨雲入酒中直欲不携筇竹杖便登南北兩高峯

穿塚刻碑成許事桓山石室果何如紛紛埋骨知無地禪智山光豈戀渠

\begin{fw}挈家人至玉泉觀魚連日臥病是日少間欣然書之\end{fw}

秋入水雲涼更白山依林薄近逾青病耽風物收微尙重過橋亭感夙經到寺鐘鳴開暝籟爭巢鴉墮散䟽翎何如策策堂堂地但許羣魚解出聽

\begin{fw}蓮士來旣作湖游復歸集吾家雜寫示之\end{fw}

未領山雲湖望寬午暄剗地換朝寒趙家園子垂垂樹猶守穠枝待我看

車行無意入雲林巖翠沈沈入畫深雁過長空影留水天衣禪語本關心

家有海棠紅壓枝先生難遣是花時閉門小病前三日晴雨相忘若底癡

海內晨星念故人逃禪不信竟沈淪近聞點茗煩雛婢却把花枝擬病身

\begin{fw}九月一夜大雷雨作\end{fw}

又是關門聽雨聲飛蚊遶鬢一鐙明何能持此除殘暑豈便相忘有早晴砧杵隣家多不寐屋廬閒地亦無情去年空費秋霖詠急電奔雷任自驚

\begin{fw}風雨後海棠盛開\end{fw}

芳訊垂垂發舊枝海棠宛宛趁花時在庭何亞三珠樹照夜眞同萬玉妃風外初開閒不覺雨中側臥看偏宜去年坐歎虛歸賞曾怪家書報我遲

\begin{fw}海棠花下短歌\end{fw}

初看有葉今有花橫枝側蘂何夭斜蜂聲作陣不捨此春光格是留吾家酣紅炫碧充簷牙當軒發色成明霞連朝風雨頗相苦粲然豈欲捐鉛華未施欄楯爲周遮對酒不可時呼茶莫敎攀折落人手高邱有女今同嗟

\begin{fw}明日更爲一篇\end{fw}

去年花發吾卜居今年花發吾檢書客言株下可埋鐵後更作花今難如物性相感原盈虛已勝論月爭蟾蜍所憂本性受戕賊遂敎一盛成終枯我今不爲強約束一任柯萼相扶疏兒童漸長亦如我歲歲花底同嬉娛

\begin{fw}湖上晚行和一浮韻\end{fw}

下窮水際上林端山色循湖一道看市外獨行原自晦城中萬瓦復初寒文身入俗時何補斂手推抨事卻難誰放光明破沈霧老夫閒爲數鐙竿

\begin{fw}我園\end{fw}

我園有一梅日往問消息高枝遲著花低枝初點的東風雖未來喜欲動顏色平生鬱此懷每可唾手得况此爲我有寢食日在側縱使閱雪霜定免受殘賊明年更作計三數當手植土壤雖不寬耙耡恃己力位置將何宜無意限南北終思我園中永不長荆棘

\begin{fw}哀師曾\end{fw}

無端江表傳消息忍哭斯人比隱淪年鬢嗟余同歲歷春游避客憶湖塵\jz{君年與余同歲今春來湖上竟留語以詩來而未見過}曰歸願竟虛前約不暝知猶戀老親私諡會當書孝潔母喪衰絰闔棺身

\begin{fw}上巳後一日挈家人兒輩出游湖上竟日\end{fw}

輕輿飽看環湖山兩隄束湖如半環臨流况際風日美從我恰得家人閒南屛午靜聞鐘語認塔爭言吾到處兒童摹擬說山川尙辦乘船呼伴侶淨慈香火明祠龕人謂禮佛多祈蠶出門已見桑葉大間以高柳垂毿毿行宮舊鏁孤山址今闢公園任人至一堂高踞萬花中置茗風廊徵往事詩碑在山覆以亭略舉沿革兒能聽果敎長此謝憂患髮鬢不染將重青人生行樂須隨分我亦年來忘喜慍全家日作湖上游一任紛紛笑充隱

\begin{fw}入山\end{fw}

螿語蟬聲作午陰輕輿已覺入山深各持湖海平時語自歛風雷定後心天下有泉供品第伏中無暑費追尋閒門慚愧居鄕弄喻此跫然是足音

\begin{fw}重九未登高晚出湖上\end{fw}

裌衣杖策與徘徊晴嶂千重爲我開林氣遠青天不暝湖光微白月初來何人相念同秋詠佳日無忘近酒杯明歲會知腰腳健卽多風雨亦登臺

\begin{fw}湖上雜述\end{fw}

波軟湖寬櫂去遲南屛斜日落松枝錢婆留有荒祠在一劍霜寒正此時

岳墓于祠晚謁回野萍無數亦花開黃塵袴褶重湖影何日看君躍馬來\jz{謂尊隱}

\begin{fw}立春日挈兒輩看梅靈峯寺\end{fw}

便持寒臘作芳晨湖崦巾車碾陌塵家食得閒休厭寂寺梅來看恰先春兒如黃犢爭翁健隖入青芝覺徑新豈獨南枝問消息塡胷生意爲輪囷

\begin{fw}爲吳缶老題其太翁如川先生象\end{fw}

過也相知慕老坡今從顴頰見峨峨卷中人物惟今夕亂後梅花在野多\jz{象作於清同治間}解與乾坤守清氣頗聞里社念霜皤谿橋一樹垂垂發石火重看歷刼過\jz{遺象爲缶老新求得者}

\begin{fw}同吷庵過拔可寓齋戲作\end{fw}

居隣步可到風逐汗衣乾閉戶臥經日從君乞小餐蟬聲晚林遠人語暑庭寒一飽忘餘事吾曹幸自寬

\begin{fw}北海\end{fw}

何言狀北海一水似南中過櫂花光亂藏山石氣同數來忘禁籞宴坐得林風獨憶江鄕遠遲歸媿雁鴻

\begin{fw}中秋疊韻寄晦聞\end{fw}

壇園花樹在誰邊秋望方疑雁已旋往歲賦詩同此日故人言别又三年索居賴可安環堵莊語無因擾醉筵爲問陽阿善歌舞正應圓魄照清絃

\begin{fw}戲投懷藪\end{fw}

月出高空助夕涼故人相對坐風廊新泥堂上成溝畫半炧鐙光接混茫子弟兩家酣睡美文書終日署名忙扇雲汗雨炎蒸地莫笑推排各老蒼

\begin{fw}招集晦聞諸公於寓齋客去成此詩\end{fw}

難逢雲日過清明草具杯槃聚友生縞紵四方期不負鋒車一夜語猶驚\jz{晦聞昨附金陵宵車來過蘇州車幾涉險}北來應喜嘗南味醒語何堪敵囈聲客散青天見明月此時翻惜客先行

\begin{fw}晚發京師\end{fw}

北來已悔失花時南去猶能和竹枝今倚曉風聞鼓角不留殘夢在京師九關虎豹曾何恤一雁江湖有所思歸語家人應嗢噱篋書無恙面增肥

\begin{fw}自京師南歸達漢口\end{fw}

二千四百十六里\jz{車中紀里如此}塗軌無驚竟得前空橐粗能作歸計橫戈幸不擾宵眠家人相念猶天末候吏何勞問道邊臥誦舊詩還蹶起危行一往賴天全\jz{末語爲庚戌居武昌寄吳中家人詩漢口舟中夜望武昌辛亥秋去此今十五年矣}

\begin{fw}一月十日雪憶去冬南歸江舶遇大雪\end{fw}

浮江初見雪歸里便經年皓皎曾無愧間關亦自賢橫天寒可忍蹙浪勢無前小閣今匡坐惟聽竹樹邊

\begin{fw}將挈家去杭感賦\end{fw}

悠悠鄕里未能謀去去江湖尙獨憂凡楚存亡任天意曹劉吞攫謝時流四年兩度逢多難十口全家住一樓\jz{甲子秋余避地時有八口全家住一樓之句近兩歲中又舉二男一女}檻竹庭梅待歸視豈能海角久淹留

\begin{fw}讀秋岳見示詩卷寫上一篇\end{fw}

君詩獨好潔機杼能自新亦若居京師不受京師塵湖遊今過我同此浩蕩春甲乙一卷詩相示尤殷勤掊擊無可言胷臆爲崚峋此豈文字間往往埋天眞喪亂嗟今時有語皆愁呻江山可助君快意今爲陳紛紛飾詞貌何如岸角巾

\begin{fw}歲晚來視笙伯仲姊留五日别去\end{fw}

三月别去遂涉秋七月同作青溪游何圖歲晚得再見况復五日能羈留一年蹤跡我如此念兄與姊行白頭南居各有家室樂北望曾爲京師憂\jz{今年秋冬余阻兵於京師}卽今坐語還太息寄書不達空傳郵二親若在念行役豈許南北同飄流將歸復憶去年事八口避地租隣樓雙梅照座送除夕餠餌晨餽午不休兩家兒女漸長大新嬰入抱還咿嚘天寒欲雪今曰歸車行惘惘資輪輈長塗幸未塞荆棘隔夜先屬加衣裘家人候門簇鐙火宵談直可忘更籌鬢鬚日老離日多長歌未足銷繁憂

\begin{fw}歸車\end{fw}

久晴雲作赤將暝樹先昏倦客仍行役歸車過水村大星看出沒羣動失飛翻暘雨皆關性嗟今未足論

\begin{fw}十二月十五日月色甚皓過拔可寓園\end{fw}

庭院蕭然十畝強我來月正滿長廊十年不出猶今日八口安歸賸此鄕記取築亭題碩果卻憐居市伴荒傖平生相勵無慚負舉債先求裂券償\jz{是夜所語及此}

\begin{fw}五十二歲生日天雨家人言余往歲生日亦多雨意以相慰感成一詩\end{fw}

五十二生日休論陰與晴本來成幻住早悔試啼聲蒙被原無勇持觴可壓驚崢嶸看歲月催遣一詩成

\begin{fw}前十五年余與拔可客武昌同主建江橋以通漢口垂成而中輟近聞復有持此議者叠韻寄拔可\end{fw}

往事同君感百端江橋一議衆爭看架舟未信將通險利涉先期不病寒何意吾言今尙用卽平蜀道後非難明明黃鵠磯邊語慚愧鮎魚上竹竿

\begin{fw}丁卯元旦雨疊韻\end{fw}

元旦初除臘高天忽吝晴積寒蘇地力奔瀑到簷聲海近魚龍眩軍過鵝鴨驚昨宵非守歲百感夢難成

\begin{fw}人日挈家至上海是日吷庵秋岳有倡和詩依韻自述\end{fw}

避風豈欲傚爰居車過村閭感歲初結習未忘三宿戀繁憂稍遣一朝除獨憎眼疾多蟬翳尙守書囊半蠹餘早悔言兵晚聞道知無困辱到橫渠

\begin{fw}題齋壁\end{fw}

冬暖無霜萬瓦乾全家鐙火共高寒中天海月橫空見捎地風篁靜夜看往日歌呼今尙念少時哀樂已無端不須更着尋常語薄醉初能夢寐安

\begin{fw}晚歸\end{fw}

晚木風多秋作晴平原能見海雲生萬書堆几蒙頭坐今許長空着眼明

\begin{fw}簡一浮疊韻\end{fw}

江近欲無岸霧深終得晴兵塵今日語礮石昨宵聲魚爛誰相恤鷗閒了不驚黃巾儻知避我獨念康成

\begin{fw}哀缶翁四十韻\end{fw}

病以不全全我昔贈翁句翁病今不甦言哀忍無語思翁蘇州時一見貽我詩署姓誤作褚明日來致詞言尋桂河坊時登癖斯堂書畫徇夙好每見開橐囊時我尤嗜印舉贈不自惜佳石爲奏刀竵扁布朱白思君武昌時遠來將北征同車樂旦夕遂作京師行朱邸好結客翁復稱我賢置酒張君\jz{弁羣}家日爲開廣筵吾時方跌蕩伎飲且招翁雜賓縱滿座翁但謝以聾思君杭州時十載曾四來梵誦禮所親頭白猶深哀又嘗居我家婦稚同笑譁撫頂譽我兒往往龍驎誇挂壁有翁畫自謂後弗如治具稱廚饌肉脯佐笋蔬最後在今夏翁能遲我歸涼堂俯湖風論畫語通微銅象近在龕眉鬢依龍泓嘗謂千載後誰更知兩翁海上十餘年事事猶在眼傷哉樓居人再見陵谷變賃居昔對街言别後還里别久僅三月書來輒累紙爲言近過從惟有朱\jz{漚尹先生}与王\jz{一亭}合君得爲三庶可相頡頏避地歲一來甲子迄丁卯翁憐我日貧我喜翁不老聚語雜悲歡往往如家人詩成必寫示琢句彌生新或許我點竄或許我塗抹濱病付一編猶強我加墨哭翁今有淚和墨紙爲溼吾叔近告徂痛與翁同日我詩告翁靈難盡翁生平翁自有不朽後當書其銘

\begin{fw}挽嗇翁世丈\end{fw}

記我初識翁翁年五十二快若平生懽相勖四方志讀書在有用但戒不習吏馮依皆自爲侃侃日論議本州苦經營間亦及國事我懷翁所知許不爲物累乙卯翁在都南歸驚宵寐呼鐙起寫詩晨發獨相示合眼此十年愍亂若五季翁歸不再出邱壑頗自恣五月翁生朝觥祝我一至猶言養高德庶可息傾詖笑慰雜嘲弄爰居鐘鼓避國非不可爲貨惡棄於地别來甫聞病凶問遽愕眙論定在闔棺詩成未收淚有才奮早施垂暮簡旁嗜此語可長生慟哉今永閟

\begin{fw}磵秋屬題校史圖久未以應歲晚乃賦寄之\end{fw}

校史圖成萬軸空冀州點勘義無同人間重視歸方本法乳曾傳惜抱翁

攬卷無端成死生知君涕淚益縱橫徐君今付龍門筆莫使荆卿見姓名

\begin{fw}敬鄕樓圖卷爲黃溯初題\end{fw}

永嘉之學首季宣寔從道潔源伊川禮樂制度求措用所學下足安元元文節\jz{陳止齋}忠定\jz{葉水心}相後先鄕守其學垂千年仲容晚出我得見誦禮法墨曾周旋君今築樓求遺編不獨數此二三賢發幽啟翳窮里壤其人縱往尊其言佳刻去歲墮我前紙墨不朽君能傳此樓此圖無間然何愁陵谷爲變遷

\begin{fw}題鶴亭爲其母周太夫人追繪稽山負土圖\jz{光緒甲辰太夫人葬其父季貺先生於吾鄕之木冠山後十年甲子鶴亭始謁其外大父墓故有此圖時距太夫人棄養已三年矣}\end{fw}

留杭不安越躑躅吾深思里弄未一詣媿君登會稽君初除母服來杭吾見之一日晨渡江事迫不可覉歸來述何語誦母孝與慈今謁外祖墓慟未書豐碑言母昔營葬女孫猶得隨三年何日月母也不及知外家賸紀聞同於洪卷葹遺書遭火厄曩僅刊遺詩防護付何人謁墓難陳詞母孝君能承世變禮就衰謁墓不憚遠亦足紆君悲乙丑歲方秋逢君於京師一圖狀髣髴如母營葬時屬題詩已成不意成分馳短篇爲補綴復後三年期先壠感霜露同爲無母兒

\begin{fw}視旭東\end{fw}

泰亭山色復當前雜以山花是杜鵑谿路漲通明岸柳陌風晴㬉散車煙君知歸客應相念老厭兵塵豈自賢尙幸兩家春事盛蜂喧蝶鬧任紛然

\begin{fw}茶具垢敝家人請易以新者戲以詩止之用多字韻\end{fw}

從我三年積垢多竹壚烹汲費清波提携曾共驅馳日棄置難爲決絕歌味水但忘年力減論材莫問客能何高曾俎豆桮棬在笑看人前輦載過

\begin{fw}紀十年前漚尹先生戲贈句\end{fw}

可憐絕代諸眞長贏得江湖七字詩\jz{二語爲漚尹翁見贈}此語旁人難卒聽紅燈呼酒看花時

\begin{fw}元夕雨姬人具酒余嬾近杯杓賦詩示之並寄家人\end{fw}

屐響高衢雨腳齊眼前蹴踏已春泥經年佳節誰逾此近海東風漸壓西家食得安留婦稚鐙筵相念索餹䬾齋廚貯酒無杯杓擁髻惟敎對小妻

\begin{fw}中秋念去年與晦聞連夕看月於壇園賦寄\end{fw}

隔歲園游此夜情還期萬里共陰晴寒忘風露仍堅坐月隔樓廊作澹明興仆何堪言世變别離可慰是勞生年來烽燧看南朔秋正平分夢太平

\begin{fw}中秋客散遂早臥明日晨起漫書\end{fw}

齋廚燈火任蕭然節日家居豈自賢去歲未歸方遠客橫天有月此酣眠不祥文字憎兵甲漸老情懷謝管絃起視中庭設瓜果昨宵猶爇燭如椽

\begin{fw}壽劉三母夫人七十\end{fw}

軾若爲滂母許之堂堂宜使後人知奉親甘作逃名計結客曾聞戒飲詞白髮娛觴無量壽黃花環屋盛開枝還勞投轄招吾醉來看萊衣舞小兒

\begin{fw}近感呈炎午\end{fw}

晨羃重陰午便晴憤將持此擬平生艱危尙辦移家計躑躅猶遲度海行還信關中勞景略豈慚車後逐康成揲蓍扐筮誠相念爲語庚庚有大橫\jz{是日炎兄爲余卜出處故及之戊辰十一月二日}

\begin{fw}賃居與恕齋故居及太夷海藏樓爲近\end{fw}

恕齋無舊牓園樹換樓扉更念海藏叟未從津上歸賃居今獨近論世早全非夾道看榆柳應敎大十圍

滬西何託迹鋒鏑復橫天未曉憎隣鬨將眠見爨煙津梁誠倦矣環堵任蕭然有客勞相慰移家策萬全

\begin{fw}哀梁卓如\end{fw}

平生文字見瑰奇聞繫安危世共知牖下得終嗟委蛻囊中探智媿同時身能卅載言興廢後惜羣流失導師致疾無端醫復悞刳肝鏤腎果何爲

\begin{fw}暑夜同炎午三兄登市樓\end{fw}

欄楯凌空一再登君猶不仗手中籐雨餘仍見朦朧月宵熱微憎上下燈茗坐竟爲傭保重笙歌豈解亂離曾避人初可忘喧寂起倚茲樓最上層\jz{戊辰六月}

\begin{fw}立秋登市樓所謂樂園者賦呈炎午\end{fw}

層構盤旋久不休二翁猶作少年游樓臺天半疑無暑風露宵來正立秋久惜衢廛非我有况驚河漢向西流行窩相許求安樂暍死應憐海外洲\jz{時聞美洲酷熱有暍死者}

\begin{fw}亞子來云楚傖見念遂賦寄\end{fw}

索米長安大有人郊天楦鼓惜騏驎覇才應笑難栖越往事曾同不帝秦狂乞白衣逃鏁試願從滄海息揚塵因君寄訊陳同甫\jz{謂巢南}葉適當年誼頗親

\begin{fw}晦聞旣來海上早晚卽歸粵以長篇貽之乙丑迄今别逾三載中所璅璅或僅我二人能喻也\end{fw}

北亂吾南歸南歸亂遂繼日月何昭昭不信已三歲幸守平生誠忍下别時涕别詩語珍重勖我志能勵不客東諸侯亦再謝書幣疴疢君乃蒙得甦如換世社園秋見月書扇寄新製圓靈同一天望望墮荒裔譏詰塞道塗箋繒爲沈滯度遼君今來豈更夢幽薊聘姬疑再求敎兒憐不慧申眉復何言有女已相壻惟隘與不恭久視等夷惠衆中早嫌身求遯難得筮袖手非本懷徇人失至計縱受投璧譏愼勿懼旁睨雞癕與豕零斯時孰爲帝學荒困戶說禮失哀野祭何日從君游與人言孝弟

\begin{fw}十二月四日寫寄劉三\end{fw}

兩宵多被酒實則叢百憂方冬天不寒氣暖身爲柔况余感初度明日歲一周上壽任羣稚近乃隔本州息念平生勞當得千杯酬有勞人不恤往事誰咿嚘是以從容飲霑醉無他謀前夕從秦翁\jz{子質}座有汪\jz{閑止}吳\jz{仲言}留秦翁頗厚我尼我無遠游昨夕從惲翁\jz{季申}論詩如射侯會者十餘人頽然多白頭\jz{是日集者爲綠筠琴士懺餘子大篤友文甫通百守梅梅光英齊恕遣卐广改廬子肅豪生與季申及其子公樾}前歸渴無飲昨歸傾茗甌卽復擁衾臥鼾息何齁齁高枕本憚醒夢往不可收友人熟粵俗云酒無近喉粵行我果成當難酩酊求因之奮我勇持觥竟不休但幸醉益恭不貽儀狄羞我年五十四下視千載不忘憂果此物一日千春秋誰歟可洞曉華涇有病劉詩成卽寫寄如在黃葉樓

\begin{fw}雪晴晨起\jz{十二月二十二日}\end{fw}

霧合猶無日誰平凍雀爭樓寒留雪意砧起辨晴聲歲盡重更朔年豐浪得名從容忘噩夢但遣我詩成

\begin{fw}過儆廬\end{fw}

婦病君能不出門學書擁卷換朝昏放閒我輩原天幸歷刼年來盡佛恩詩卷狂題嫌紙短座氈久暖勝爐溫雪晴過巷無他客海市眞看靜似邨

\begin{fw}朱三兄炎午爲余致歲暮之資感爲一詩\end{fw}

朱兄志嫉俗干人不以私獨爲歲寒友取供歲暮資嗟我恥干人居貧人不知舍己時助人貧亦能自怡急難恃友生輒在空橐時攫金不署券往往顙有泚治進亂亦進久念相勖詞奮身果有濟殖利終不爲君昨與我言貸者楊翁貲曩曾通緩急今故能致之挹注感相恤偃蹇負所期固知平生交不以責報施託身在天壤未嘗一日饑捫腹縱得飽何能坐而嬉儉無人索債有客誦我詩\jz{此乙丑歲除詩友人胡君栗長喜誦之}以儉仍居貧日月悔不追

\begin{fw}寄晦聞\end{fw}

北亂曾同感未歸我歸今後别湖扉久知俗敝無愚智終信羣囂有是非長夜笳聲逢曉靜一春花事任風飛越中劍氣今沈歇恥著曼胡短後衣

\begin{fw}贈梅泉\end{fw}

我少婚桐城奉手隨馬\jz{通伯丈其昶}姚\jz{仲實師與叔節三丈永概}尤感徐翁知\jz{椒岑丈至南昌索余相見}忘年曾下交囊書走海上汪\jz{允宗}李\jz{曉暾}相招邀是以皖中賢能言我所遭近歲識周君鷦寄榜其巢爲政誦所先論詩冠於曹席業本非豐自致稱盈饒意仍有不屑捨擲無訾謷致害可早散志曠原忘勞僶俛五十年紀歲今如陶置膝已抱孫豈獨羣兒嬌紛紛世間事善歷窮所繇期君奮華首折義平衆囂

\begin{fw}贈伯夔\end{fw}

君家數兄弟識君惟獨後吾友張與湯譽君不去口念同列臺試四海孰誰某其名曰同年實遜文字友奉親海上居頭白君有母朝籍雖一通陔養樂長守君益肆於學近師散原叟文字如其人簡重復婉厚誦君有軼事舉措誠不苟甬人昔發難爭地奮擊棓異族求戈君獄急難出走君闇遺以金使得脫械杻我曾理此獄其事知八九祕牘曾親焚於君爲不負捉筆今書之君當頷其首今年君五十期我醉以酒未醉先成詩知能異它手觵觵誦先澤爲澤定遠久相期葆黃髮君更爲我壽

\begin{fw}炎午拉過西婦海倫巴勃試手相術漫述一首\end{fw}

西方敎論輕鬼神蕃婦近乃眩相人投之以錢試其術出手視掌承几茵鉤輈吐語譯可辨初亦刺促談生辰繼言壽可七十五骨相歲率差相均更言三婦今喪一五男得與淵明鄰繼言致富可殖產此爲一笑忘吾貧朱翁同問所答異星印在指官除眞又言平生勇任事胸次了了窺輪囷有兒旣壯可别處有婦近健無病呻四客可來釋煩慮客來期以明歲春其詞將竟復它試一球交握光\jz{字侖原王}璘置盂平視它不解謂中有光相屈伸休咎一一述先後下亦可慰無灾屯自謂其術出印度豈大圓鏡光明輪姑布子卿久不作鳶肩火色今紛陳眼中豈有酒家媼能識劉季將亡秦朱翁與我各偃蹇閉戶久矣忘海濱黃金脫手任早散白首傾蓋期如新兒曹美劣不足數但願人海藏吾身彭殤一致齊萬物何者弱草何靈椿有文在手古所誦未可執此相逡巡其說果驗亦細事奮腕不樂隨風塵異哉楊君\jz{恭甫}昨相遇捫骨望氣如引申歸爲賦詩紀名氏氏曰巴勃名海倫

\begin{fw}午倦賦禽其有類於琴操耶\end{fw}

出巢久苦不得飽居巢自信甯忍饑巢中伊啞四五子中有嫁者雄挾雌老禽蜷伏恨無力自奮欲傚初飛時巢幸得完子數生晨翔晚飼巢中兒網羅不及愛文采雉乎亦受兒童譏世間每詫希有鳥實則毛羽皆離披欲出不出木有枝已病未病心有思竹實可食鳳所私吾如大鳥其從之

\begin{fw}贈劉季平\end{fw}

海上飄鐙失歲年但驚天轉地迴旋忘憂未敢從人說握手相看識汝賢論誼直宜詩作券當筵何惜酒如川醉中歌更聞高詠第一須題寶劍篇

\begin{fw}畏暑輒書一篇時左趾方病濕\end{fw}

園葵遲開久不花庭梧已滿旁生芽今年太白說晝見何怪寒燠皆爲差五月井溢生泥沙兼旬雨腳紛如麻連朝驟熱解巾韈瘡疻被跗窮搔爬視若癬疥本無患所愁嗜鰒甘其痂古人歲厄歌龍蛇後來月蝕唫蝦蟆平生觀物自了了有兒對笑相啞啞願兒莫更如阿耶近耽野性時居家自晨及暝不踰閾一任夸父追日車

\begin{fw}正月十五日挈女兒甥輩謁仲姊殯舍\end{fw}

節日何堪聽哭聲殯廬偕謁女隨甥平時篤愛全留此往疾猶疑倖可生歸葬聞將求故里燒鐙忍更說春城楚嬃應識三閭感䩌悴行吟久不情

\begin{fw}己巳正月二十日女兒雲怡來别將赴廣西時壽氏壻方客柳州也以詩示之\end{fw}

汝先失母今爲母提挈雙雛桂管行蠻瘴竟忘非遠地海航端念是初程明知我健無衰病更抱孫歸見長成刺促此時相慰語老懷剛喜雨能晴\jz{是日雨晚歸值霽}

\begin{fw}二月十一日爲三兒安世生日賦寄家人\end{fw}

汝生在鄰巷記取五年前墮地爲吾喜留家伴母眠扶牀今解事索果不求錢能葆童心在他時許象賢

\begin{fw}近詩寄朱三兄炎午\jz{五月廿二自南京寄}\end{fw}

端策何從問卜居未焚賴有腹中書一椽託命移家後萬事安心失望餘便逐征鴻别巢燕欲忘轍鮒累池魚虛空粉碎吾曾悟持擬吾生豈不如

\begin{fw}三日雪中期炎午三兄可來聚又念雪行中之勞感成一篇以迎之\end{fw}

冬暄何意雪浮空書到猶言道可通久欲掃除塵榻待今憐辛苦驛程中孝陵下拜尋苔石淮榭相招乞藥籠\jz{時正目疾}便與安排成日課願君行李莫倥偬

\begin{fw}題缶翁山水卷前廿二年曾屬題而未果者\end{fw}

展卷重爲思往日尙疑翁未棄人間亂頭麤服一家法淡染濃皴幾點山直似草書多媚嫵自將篆勢作回還當年我語翁能喜畫訣通微古可攀

\begin{fw}壽太夷七十\end{fw}

去秋君謂春當歸將不再出還初衣從亡未遂沈海烈論壽亦是懸車時平生重君守節概在閩可洗劉氏譏日星爲紀齒縱老勝我五十知昨非當年相逢各壯盛玉貌欲解邯鄲圍贈篇相許同魯連所語成讖嗟退之救亡二策震一世寺廊往語爲累欷國如藏舟夜負走家有重器鄰爭窺君申論難我奮筆意可假手扶顚危奔騰日月今廿載賦詩壽君何及斯往懷早見鷗榭記衆譽或涉龍州師世更誦君書及詩寒泉落澗林禽飛君歸行復值生日我言儻爲傾深巵惟忠與孝可永永侍次况有君佳兒

\begin{fw}青溪傚玉溪生體\end{fw}

婑媠妝成鏡有塵旁觀淺笑復輕顰束腰欲謝長裙窄斷髮翻求峨髻新守禮東鄰猶處子擅歌南國有佳人蔣家小妹誰相訊聞道青溪正賽神

\begin{fw}泰淮放舟雜書所感得十八絕句以示同游\end{fw}

橋北橋南柳色新晚雲將歛見鱗鱗淮流不照驚鴻影但照當年縱酒人

大笠疲驢入定林一時師友詡清吟晚歸淮榭從容語多是人天未了心\jz{謂楊仁老與桂伯華師}

豪游曾共憶陳\jz{散原丈}俞\jz{觚庵}眼裏樓扉認故居留得濠堂分一水\jz{濠堂太夷所居}同看頭白感黃壚

舵樓初見許生平施叟\jz{理卿}相偁坐可驚若輩豈知家國事我曾揮手謝公卿

李生\jz{曉暾}埋骨斷知聞更慟劉君\jz{申叔}繼魏君\jz{季詞}歌哭獨容稱健在無心文字醉紅裙

容許清狂閱七年奇章能識杜樊川當時曾博雙鬟笑譽我逢人在酒邊\jz{謂恕齋}

寶華庵裏一尊同騎吏來迎問醉中散髮斜簪工捉筆故人眞有晉賢風\jz{謂匋齋}

鐵畫樓中老尙書佳兒名父許相如可憐繒繳終難免未許王尼宿露車\jz{謂張仲宅塏曾}

小字眞珠寫柱銘一時已惜累書名何期老作諸侯客盾墨猶爲列校驚

曲檻明鐙雜鏡臺舊曾扶醉此徘徊傳呼隔岸河廳散未許輕輿過巷來

列炬樓船夜若城鞾刀帕首氣縱橫今看早燕初鸎裏仍有喑嗚叱咤聲

一雙笛婢楊公濟海上爭傳事有無不屑流連竟歸去心腸木石是狂奴

舊時潘薛半飄零擁髻當筵語忍聽終悔流傳有篇什千金漫許贖娉婷

輕舸沿緣打水圍卷簾爭出泥遲歸廿年往事橫胸臆座上誰知白袷衣

東關一曲尙通潮舊院難尋武定橋誰遣笛聲工入破醉懷難遣是今宵

觀河面皺本無端自懺風花語亦難行盡秦淮三四里已無人倚畫闌干

隔月來游倦可知晚涼仍是換衣時前塵未信眞如夢寫入青溪本事詩

未許青溪屬蔣家更無頓老擅琵琶吾曹但可求清境不乞南朝岕片茶

\begin{fw}杭居被焚歸視感賦\jz{三月二十六日}\end{fw}

宅已火中盡吾疑亂後歸花枯仍倚壁書燼更無衣墮甑難迴顧焚巢可悟幾重哀同厄者鄰巷過人稀

舉家爭告語驚竄得生全成毁非關我艱危倍視前來看有焦土自念已華顛歷歷垣牆在曾居十二年

車行疲晝夜夢囈歷川原何語勉羣稚終期大我門浮家嗟昔慣先澤幸今存\jz{先像家乘與先君遺詩稿及手寫張江陵集幸早移藏上海}再創中興策難忘故舊言\jz{炎午儆廬拔可旭東笙哥景弟皆書問相繼爲籌後計}

焚時吾客外歲厄早爲憂兀兀生何患熊熊影尙留事過猶戒愼躬責在愆尤蒙難艱貞語時哉忍自謀

椽瓦燒無賸甁罌掘尙新挈行殘硯在來問故人親藜藿求生計蓬蒿惜此身本來無一物今掃甑中塵

\begin{fw}王庵種花樹詩\jz{幷序}\end{fw}

庵旣新葺桑竹翳如一杏二桃強稱花事遂於庭場雜治花樹於塘上植芙蓉六十本自晨迄午植援甫竣雨乃如注園丁忻然稱可皆活春開秋落時序可知不出戶庭以花爲歷興有所寄各繫以詩

梧桐非異質材乃中作琴不如全其天斧斤莫相尋留此徑丈枝自有太古音占土本不多十載期成林\jz{梧桐}

庵垣築新就垣旁植以槐其陰下覆屋當暑無炎埃但恨道旁人不得隨徘徊蔭暍負本心此槐誠惜哉\jz{槐}

七年看海棠城中皆遜我新叢今遠致明歲開萬朵溫公營獨樂園扉不敎鏁慇懃告遊者攀折莫爲禍\jz{海棠}

南山多此花春發如屛風時疑五色雲幻入晴嵐中二株何娟娟定逐啼禽紅春禽啼不息花儻能相同\jz{杜鵑}

吳中昔賃宅花光照屋椽江船爲賦詩魂隔聽雨眠漂搖尙倔彊慨詠復海邊日月何崢嶸睠兹十四年\jz{辛夷庚戌武昌賦詩寄吳中家人有近聞春訊傳吳語已有花光照屋椽及那料西來聽雨眠之句其時居蘇州金師巷樓前此花亭亭如蓋壬子居上海復爲此花賦詩所謂倔強能如我漂搖尙說春也}

其葉擬冬青歲寒能不凋方春盛跗萼白者如瓊瑤偉哉霜雪姿不作茗荈嬌誰書綠莊嚴請以榜吾寮\jz{山茶}

吾無隱乎爾一語遂通禪儒門有公案誰似涪翁賢垂垂風露中秋氣來無邊佛香日在几與汝爲因緣\jz{桂}

故屋久易主故根在南昌先君曾歎息繁花照虛廊涉想時來前紅紫相扶將遮眼與落手魂夢終難忘\jz{紫荆}

庭榴歲不實花開幾至冬其葉復冬綠偃仰如園松移根費灌溉不求旦夕供高秋見爛熳輔以木芙蓉\jz{石榴}

芙蓉性喜水種日復雨中方塘步四周期汝爲秋容園人昨語我伴者惟霜楓年來已止酒醉頰難爲紅\jz{芙蓉}

籬隙不種花植此有苦心三年數柯葉應作十畝陰春來摘新芽味可充釜鬵薰蕕愼自擇功罪惟汝任\jz{椿}

託命依牆籬花開難見午幸汝得佳名使識田舍苦昨栽疑不活蔓發已可數秋曉待來看明媚此環堵\jz{牽牛}

\begin{fw}王庵雜詩\jz{旬日數至晴雨異時雜書不加編次}\end{fw}

半道江經過獨尋西馬塍村收湖外雨坐倚佛前鐙簑笠行將慣泥塗辱未能任擔禪版去粥飯卽爲僧

園荒先闢徑廊盡始登堂我自忘喧寂客休論短長麥風占礎潤林日射門黃晴雨關何事谿山是此鄕

朝朝禮名塔林壑見西南不信閱多刼猶能賸此庵雲門爲眷屬良士是奇男\jz{王仲瞿與其妻金雲門昔曾居此故今名王庵良士亦仲瞿字也}遠矣乾嘉事今供野老談

易長不速朽平生種樹心已遲春後計還向雨中尋可辨非膏壤惟期有綠陰腁胝竟何補使我爲沈吟

移屋成新構支牆有故碑應知彭尺木曾記放生池\jz{庵有彭二林爲清甯寺作放生池碑}馴鵲分廚飯游魚避釣絲眼前觀物意幻住復何辭

近地多墟墓誰持麥飯來禮衰虛設祭緣盡本無哀穿冢隣將卜銘幽石可開分明數行樹待我手親栽

早辦青山臥今爲白髮人田居猶苦逸壁立未全貧元亮知農事林宗戴角巾任敎入圖畫還恐損天眞

午出歸曛黑車行亦告勞時來驚里巷好事及兒曹得雨花都活編籬筍漸高旁人莫相恤仲蔚樂蓬蒿

\begin{fw}病中雜書所感\end{fw}

一國皆不知獨知吾其危此語出箕子世乃稱明夷殷亡周復興易置如奕碁家國若強預頌者謂帝師我知箕子志洪範皆支詞

不雨已憚出衝雨驅短轅疲馬戀箕豆鞭策多瘢痕非此不能御騶卒日有言銜勒久未除皮骨今空存障泥莫自惜人與馬何恩

\begin{fw}病中憶范丈無錯昔年所見語者爲詩紀之幷傚其體\end{fw}

昔聞范翁言病臥憎羣兒揮斥不忍去復逐羣兒嬉其病亦霍然此理不可推余時方壯健過耳忘若遺攖疾今伏枕婦稚時相隨大兒頗解事戶扃不內窺小兒素跳踉屛息若有思婉孌二女子侍我日就醫登車輒敷茵聞呻屢顰眉次女本多病新歲強自支視我日就瘥家人始嘻嘻骨肉晚益親樂苦心共知病中我善恚念此爲解頤涉想及范翁其語通天倪九原不可作太息詩成時

\begin{fw}哀閩二翁\end{fw}

\begin{fw}　魏季渚\jz{瀚}\end{fw}

前遊墟墓間相地曾逢翁言自海外歸惜我前不從更言舉兩男近樂與我同壯健猶當年失喜無衰容感翁昔論我所囿在國中何如縱所之政俗窮西東甘爲作喉舌盡攬西方風蹉跎失遠游别矣如燕鴻哀耗隔歲知未一臨幽宮晚聞嗜博戲喧笑聲隆隆此爲耄歲事不用哀材雄

\begin{fw}　王又點\jz{允皙}\end{fw}

碧栖號詞客其詩寔勝人霜天老病句傳誦尤絕倫平生近閨襜解與蛾眉親千金藥自知胡爲殉其身遂使跅弛才僅與周姜隣昨方語葉\jz{遐广}夏\jz{劍丞}閩海懍潛鱗更念京師遊裙屐同笑嚬畏廬騰雅謔飾美羣疑眞君亦不謂忤隅坐言芳塵歸人後得見友好惟稚辛夕飲遂大漸此亦今幸民

\begin{fw}心白索詩二十餘年之交誼非一詩所能櫽括述其初心商其晚計遂成一篇\end{fw}

降伏皈依感一時晚能梵夾日相隨學書不進非諧俗遯世將歸尙乞醫\jz{君將於常熟謀施醫事}早爲貧兒求寶積但安羣季慰輖飢低回二十年來事強半平交未盡知

\begin{fw}賦謝季平三兄贈蠶豆\end{fw}

僦舍新移靜若村故人傍夕走呼門芳甘入饌逢蠶訊爛漫堆盤挾酒痕再摘世方憂抱蔓相煎人不恤同根敢持憤語酬佳貺恕醉田家老瓦盆

\begin{fw}立春日挈兒看梅靈峰寺\end{fw}

便持寒臘作芳辰湖崦巾車轉陌塵家食得閒翻厭寂寺梅來看恰先春兒如奔犢爭翁健隖入青芝覺徑新豈獨南枝問消息塡胸生意爲輪囷\jz{錄者按此首重復}

\begin{fw}和伯嚴\end{fw}

南都山水鬱嵬隈公臥詩軒捲碧苔小病逢秋喜清健西風來雁記初纔世方多難愁霜鬢日飲無何迫酒杯我語李生同絕倒不談王霸卽人才

\begin{fw}十二月二十日雪後車行竟日達都愾往撫今雜述以詩\end{fw}

戒行已三日朝旭始展軨豈憚風雨寒塗軌愼所經來時道方阻江渚求歸舲攀登蹈百險中夜爲長醒乘車此西邁餘悸猶竛竮兀坐悟通塞照眼吳山青

遠林若無枝遠水若不流昔聞述畫理今以證兩眸浩浩一白間時將及常州密雲曾連朝凍合遂久留垣屋盡加厚又若新堊修吾目本久病眩視輒小休人多詫壯觀顧盻忘所憂吾疑泛雲海泛泛惟虛舟

玉山能照人茲行恣奇觀恨無巖下松相倚成高寒危厓已砥平曲陂皆坎完行者果投足或可忘諸難中原本浩蕩尤感伊洛間有如香山勝乃有八節灘夷險終如何莫作雪後看

栖霞富紅葉霜後多來游今過作雪嶺崖岫皆難求積素蕩羣垢叢卉森琳璆更無赤與蒼點染成高秋意欲縱遐矚此行猶難留昔聞俞翁\jz{恪士}言淨土貴自修步步土俱淨何用師比邱亡友久宿草猛憶念不休

平生見雪多戊戌歲將除山行越二日從者舁以輿下坂每疾走踏冰復上趨省親不中懾失道曾自愚其時已宵半四望無室廬有燈有犬聲初謂卽荒墟力行始得達一叟惟噤吁出錢乞爲導久之逢村閭孤兒今頭白何以不安居

新昌有二園臨河鬱相望曰熊氏漆氏每雪登其堂吾時方少年策騎結客場敝裘不求煖志意何飛揚城南雪櫂圖甲乙題其旁今念同游人一一皆凋喪我今幸不死我髮亦已蒼不知此時雪曾否同新昌

辛丑祀竈日我從東鄕\jz{江西縣名}還母病正念兒計程兩日間道行忽大雪嶺峻不可攀急投逆旅宿未暝先掩關葦薄列三榻壁鐙照凍顏一夜雪不止一夜寒不眠晨興見隣榻陳尸方赫然始知昨新死棺遠故待遷趣行忍寒冽見母不敢言此亦雪中事追憶今何年

燕市我七至見雪惟四冬又一至津上卽逐南飛鴻信知北土異不與南中同祠祭陳后山故人多相逢或從故人飲談讌能從容判别半死生今復各西東吾詩不能寄江上多北風經過江南地今日在雪中

西湖十年住飽看湖上雪船路通淺澌柴蘆見森列積冰煖始開其聲如帛裂重煙化爲雲山崦各明滅遠岸頓使近其景尤奇絕近歲多去鄕惟湖不忍别惜哉今未歸欲辨言辭拙更念紅桕莊叢篁凍欲折忍飢甘里居此樂豈盜竊

此行無聚散兩地皆寄孥嫁女未遠行伴母同晨晡尙言别離苦自笑非丈夫到都日未旰積冰斷衢塗轍困多費時入夜猶馳驅平生念行役樂與冰雪娛到家就爐火奇溫在肌膚得煖念奇寒誰歟共崎嶇

\begin{fw}積雨樓居戲示鉢文特璋\end{fw}

雨聲雨聲挾以風一樓坐臥如舟中直疑江海走巨舶隣有機軸何隆隆牽船岸上古所喜今我不樂誠顓蒙衝泥陷淖者誰子踞樓俯視誇三翁狹扉起鐍避急溜縱談大笑惟乖龍\jz{是日談江西鐵柱宮事}龍方行雨憤相懟亦翻奔瀑傾長空家人所苦在蚊蚤中夜數起爲憂恫晝潛宵嘬蜰更黠不如飛躍容此蟲沈霖蘊溼始致此驅搏可恃惟人功明朝雨止日不出枕簟安夢宜朦朧驕陽果使灼大地更恐喘汗相忡忡年來袖手罷篙檝虛舟一任隨西東此時風雨聽橫集卻勝軟語依吳篷平生苦樂分若此我與二子今相逢

\begin{fw}暑夜聽琵琶因感昔游且念亡友愴然作\end{fw}

古樂一變爲新聲久無琴瑟琵琶箏琵琶入唐已亂雅今亦寂寂絃難鳴忽雷大小昔曾見\jz{唐琵琶大小忽雷昔爲聚卿所藏余曾見之}委棄幸未歸榛荆縱張伎舞盛歌拍誰遮半面檀槽橫暑雨苦溽秋未成高樓鐙火誇初晴市塵倦踏避曹偶不意入耳如聞鸎發喉掩抑雜激楚心苦不樂移吾情撥搊卽使失眞意視此已足娛生平近聽列座悔不遠強欲捨去終未行聲從何來出四絃一時惘惘無它營盛娥老嫁斷消息\jz{余二十年前識盛氏姬一時許爲琵琶聖手}沈侯\jz{謂小沂沈五叔兆祉曩曾戲語余跳盪少年當爲盛姬敗道今沈謝世亦十年矣}雅謔旁人驚阿宗誤我哀一語李翁晚亦言林生\jz{此謂林宗孟長民林與盛最暱阿宗者盛之稱林也}穨鬟不理淚承睫每爲再奏傾杯觥曩時綠鬢今白髮後有繼者名難爭此夕愀然聞此樂沈林逝矣歲歷更才人不遇世尙惜旁羅繁猥皆荒傖錢塘江上櫂女迎有女能彈目苦盲\jz{錢塘船伎錢鬘珠之姊工琵琶而目病盲知余嗜此在江樓曾爲彈半夜至月出始罷}眼中海曲卅年事更比錢塘江月明

\begin{fw}樓居雜興\end{fw}

賃居何待卜得地卽全生再徙依前宅\jz{今居之隣屋爲甲子乙丑避地時所居}一樓開曉晴眼明書獨校\jz{近校新唐書久未竣}世亂夢多驚懷抱無端惡兒童辨市聲

擾攘冬春際提攜子女來子愚多廢學女病竟長埋\jz{二月二十八日稚女以病殤}制淚無他語思親有積哀\jz{丁未十一月九日先大夫棄養時亦在滬}誰知不祥地使我久徘徊

召亂先亡禮驅農半習兵此憂吾早痛何日野能耕轉粟空千里馳書下十城無糧終自潰消息本分明

聚書原塞屋巾笥惜無多恃此伴哀樂直如逃澗阿雞蟲論王霸蛟鱷狎風波鴉背斜陽盡堂堂一日過

昔盛成衰落淮游感去年其時逢麥熟猶勝苦烽煙葺屋謀兵後無蔬進客前楚州城下水今阻過江船

吉安生長地吾叔近重經書不言衰狀兄猶守祖庭寄錢供半飽嗜酒念長醒避地仍家食非同聚散萍\jz{久不得三叔消息近知已往吉安依大兄惟弟姪尙留南昌}

哀郢懷湘事誰招屈宋魂術窮終自賊世狹竟難存閭巷多宵遯耰鋤竟日喧脫身來海上愁對各無言

世皆譏毒月俗尙守吾鄕雄佩兒求艾虛庭婦炷香淮游今獨念歲隔事難忘生聚謀難用荒區換戰場

\begin{fw}哀李觕巢丈\jz{宗言}\end{fw}

丈丁巳卒於京師後十三年庚午寒臘檢是年所爲詩有哀閩二翁者蓋爲魏君季渚王君又點所作追念往歲追悼丈之五言二篇竟難記憶然丈之風義行誼不能不見於吾詩補成一篇得三十四韻

江西四知府丈與徐\jz{元霖}趙\jz{于密}王\jz{以慜}忘年論文字延譽親老蒼我翁與丈游徐叟同激昂趙王時稍後相視若輩行丈旣守廣信招客寶井堂扁舟我未赴一諾今豈忘抱關強自卑將母有不遑其時數相見丈正留南昌後來走海上日坐屋下廊喜我能造門藥裹曾分將一朝聚京師驅車歷城坊雜坐隨家人願言日在旁貧官食不給囊金資其糧憐我已抱子甘語酸中腸一朝憤南歸丈獨稱其剛湖山恐坐嬾出處恐無商奈何不繼見一棺葬其鄕爲詩曾哭之行誼哀無量距此越十年宙合仍茫茫近聞最小女將嫁縫裘裳感丈書生平所師聞四方偶齋\jz{寶竹坡先生}與芳郭\jz{袁忠節公}如蘇從歐陽故山欲歸臥爲園題曰償此志縱未副先業云久荒流傳支社詩豈播文酒場在家曰孝弟在政曰循良恆行不可掩作傳疇能詳母爲文肅女\jz{沈文肅公}厥妃氏曰黃所生有六子大都工文章吾詩語止此爲時卅載強幸不負故人百變安故常約堂丈猶子今同旅江鄕如丈能厚我投分容其狂詩成一寫示願共愼篋藏

\begin{fw}衆異崑三招余聚飲余以迫睡先歸道中賦謝幷記一時興趣\end{fw}

早歸耽睡慰閒身鬨笑鐙筵失主賓馳道不離天上月小車何讓洛中人移家吳下園商借穿塚要離碣信眞俊語詞人工點染有無執著兩隨因

\begin{fw}拔可近游燕子磯有詩依韻奉和\end{fw}

巖洞緣江達懸岸誰其游者趾垂半此城燕子果飛來靖難當年曾召亂日與山靈洗謗名大江潮落復潮生岸人失色舟人險歷歷篙痕石隙行

\begin{fw}金陵赴友人之招\end{fw}

散㘘東關晚日明赤闌橋下暗潮生高城自束淮流勢舊院猶傳箏笛聲對酒不懽隣舫笑當筵無語燭花橫玉京入道橫波嫁留作人間百感幷

\begin{fw}臥疾\end{fw}

臥疾閉關遂竟日急風吹雨欲沈江潦深應沒縱橫徑林翠時窺窈窕窓濕蟻來尋看一一幽禽相對語雙雙翛然隱几忘天籟獨信繁憂未可降

\begin{fw}雨中友人招讌不赴卻詠一首\end{fw}

隱隱高衢聞屐聲已知秋雨未能晴司觴應避驚鴻影列炬遙知雜騎行大海風濤方震撼小樓絃索益淒清不如天際眞人想歷刼風花悟此生

\begin{fw}兒子章世今歲年十五於其生日以詩勉之\end{fw}

四十始生汝汝今十五年誦書能上口我意爲忻然愼哉幸自擇人有賢不賢

汝今所作書先愼後頗肆作書當主敬古人必如是兒毋傚翁書老作跁跒字

汝近多放言爲我所屢訶兒時不自斂少壯將如何士憎兹多口垂戒思孟軻

搬演出諸人幻戲原尋常汝今喜往觀雜坐閉廣場小則可致疾念及爲旁皇

兩姊爲女兄弟妹皆後汝門戶將以託善惡宜自取我生無弟兄汝應感此語

日夕在我側猶覺多荒嬉求學走四方茫昧後難知第一當猛省如在我側時

人生小天地惟有父母恩應推父母心敬愛以自存家庭相雍和其樂在一門

飲食本細事口腹莫浪求過飽易致病我心屢爲憂吾今幸少病此故汝知不

\begin{fw}賦答衆異\end{fw}

故人詩札到江碕慰我猶言不出宜身健恥言貧是病機忘時笑退成詩鴟夷生計千金少龍漢光陰歷刼知袖卻量江手無用吾曹等是一絇絲

\begin{fw}宗孟旣及於難余以聯輓之詞曰論事非有可爲君歸地下應橫涕殺身誠所不幸我媿兵間再脫身後四年浚南與余相遇於金陵出其哀宗孟詩三篇見示爲詞頗痛遂追輓亦成二篇\end{fw}

曾怪豐髯易短髭同哀時會孰爲之單車出塞終何濟雙栝題門後可悲漫信將軍能好客更無弟子爲輿尸廿年前與論官制叱咤旁驚奮臂時\jz{庚戌在都君與余論官制余主張兩級甚力君謂不可其詞色甚忿座客駭然旣而君乃折服後鈍初曾語余草地方官制當從君議宗孟亦見從矣}

消息當年阻戰塵日聞季武惜佳人\jz{林君季武日與余商君得脫歸之策且時誦佳人叱利之詩至今猶爲感悼}九原應慟難謀國四壁尤憐不諱貧爝火功名猶歷歷浮雲出處本陳陳楚些先後哀亡友海內今仍有病呻\jz{末謂鈍初}

\begin{fw}聞蟬\end{fw}

王城豈無暑清曉卽聞蟬科跣難從衆飛栖本自然厭聽人習戰傳語虜窺邊桑土將何計吾慙褦襶前

\begin{fw}小病憤書\end{fw}

十日幾無一日安心知內熱雜中寒蓬頭椎髻將何慰\jz{時家人兩以書來問病}虛室單床賴自寬往迹可忘同一埽此生早分閱諸艱祝宗祈死非吾事蔬筍誰從飽晚餐

\begin{fw}俶仁得陳簠齋所拓秦瓦量精本屬賦長什遂爲五言十八韻\end{fw}

漢律文多佚秦量字可捫篆留刊詔令制重壹黎元徙木民曾犯銷兵蹟豈存亡胡先應讖二世已無孫創業兼倂始求仙汗漫言斬蛇開繼祚指鹿誤童昏不盡咸陽火徒傳枳道恩用知成一統未足守中原功任銘金石人猶寶瓦甗臣斯詞具獄微管器同論鶉首天何醉狐鳴衆早奔風難歌駟鐵牧已失乘騵偶語人重足詩書刼一燔荒荒看世換往往爲聲吞異代何餘此當年用未繁斗升藏怨黷雕鏤役精魂愾息非嬴氏奇寒落海村千金知不易呂覽等懸門

\begin{fw}得衆異書卻寄\end{fw}

度遼君久去安越我無言人海勞緘札書巢念閉門五年秋似昔百病舌猶存黑石臨星浦潮聲聽吐吞\jz{得衆異大連黑石礁所發書與衆異不見亦五年矣}

\begin{fw}過吷庵所居折花以歸\end{fw}

淺白穠黃作幾叢春餘雨挾夏初風園林得地居何遠蜂蝶隨車日正中便欲移栽無隙土尙容分餉到隣翁歸來火急求甁甗几榻凝塵一掃空

\begin{fw}方暑多雨時拔可夢旦崑山衆異往游天童雪竇二寺\end{fw}

山游言避熱先犯海程寒雪竇竹當見天童雨不漧巖栖思靜地鄕味戒行餐腰脚吾慙健詩成定許看

\begin{fw}妙高臺和衆異韻\end{fw}

我聞妙高臺縹緲化人城溪山占幽勝緇素皆經營幻住果何有風松不留聲土木誇一時有情徇無情游者能一至虎迹迷縱橫崇樓俯廢塔過視誰重輕山中人有言從浴疑華清雲窗與霧閣豈得山靈盟

\begin{fw}拔可旣和衆異妙高臺詩更依韻一首\end{fw}

名山古游幸列騎如堅城况近樓山村何殊夾馬營紅粧從娥娥刀劍應無聲荒塔寺重堊輝映若有情念以緘白骨涕淚當縱橫美人果何物休令黃屋輕前魚亦可泣濁井不再清妙高留此臺忍見釵鈿盟

\begin{fw}聞奉天陷於日本\jz{九月十九夜}\end{fw}

內訌召外侮累載垂至今江淮痛告災遼瀋警來侵奉天一朝陷忍復聞吉林毁車殪渠帥其勢本駸駸蒙面以事仇忘恥同無心此卽足致亡不待鳴笳音久疑非我土旌旆何沈沈叩關號習戰坐視皆爲瘖嗟哉萬寶山後將不可尋亦如皇姑屯世不哀作霖燕雲十六州委棄歸遼金其源在五季此辱同呻吟

\begin{fw}辛未秋用儆廬韻賦悉蟀\end{fw}

東北驟易候蟲音聞益苦其音有何詞瓜豆訴分剖衆命託刀砧連戍罷笳鼓不敢撓厥鋒直疑貓似虎歛退復追躡奮張更雄武二十二萬衆難與萬爲賭勝者鳴其羣弱者戰其股天生德於予漸臺笑新莽悲哉炫秋聲久矣成奕譜蠶食葉空筐鯨吞海無土矯行烈士烈坐歎腐儒腐蟲兮汝求生天意定能許

\begin{fw}讀史未竟聞客語有述\jz{三月三日}\end{fw}

上巳元宵集一時獨傳消息累人疑使兵環守原無罪舉座相驚旤可知澣日眞容就湯沐流言或見卽誅夷班書久爲分人表晁錯類如蕭望之

\begin{fw}暑中將謀徙居姬人復病賦此自遣\end{fw}

占室何期十笏寬憂時僅祝一隅安雨聲歇處啼螿喜海氣生時倦鳥寒漸老敢言家有託非秋已覺感無端病姬付與歸心法火宅清涼本大難

\begin{fw}暑中以詩狀之不加詮次\end{fw}

隣廧若危厓其色著深赭初陽換晚日吾廬影其下時疑山中居兀兀過長夏惜無千尺瀑傾瀉補檐罅毛髮爲灑然不媿山居者

門前有隙地雜坐婦與嬰日夕豈待招攢次若露營比隣吾漸熟不問知姓名人雖產贛粵語皆辨其聲衆意在追涼久久熱自生

樓下置涼榻兒輩爭就眠一朝暴風雨徙榻相後先兒豈有愛憎時自有變遷爪觜方紛拏壹意求安便我意難語兒拊枕爲茫然

道旁見車人街陌事糞除其來日未明行者不在衢市風樂晏起轔轔聞過車擾夢或不安掃地非所須晨興吾嗜早每見多感吁人我果皆淨酷熱還太虛

履穿露其踵古語狀人貧近人性好惰靸履何紛綸我國豈貧國習俗遂移人毋勞削趾適直與徒跣親何聲自巷來躡屐傚東隣視此猶勝彼世變何陳陳

\begin{fw}爲衆異題王阮亭少年晚年手寫詩卷\end{fw}

早得虞山一叟知落箋堂集始編時稹貽溪上分名字卻惜當年未入詩\jz{文簡名字蓋以司空圖隱稹貽溪上之義見鮑鉁稗勺故拈出之}

正宗絕代託銷魂申叔曾爲語辨寃獨負遺民哀小雅顧亭林及杜茶村

彈指華嚴示現中樓臺七寶本玲瓏嶺南江左爭壇坫合讓新城絕句工

我藏手校漆園書持擬君家妙跡如隻字寸縑應篋衍論年四十九年餘

\begin{fw}拔可連得伊墨卿書幅皆舉以走視且云將榜所居曰墨巢爲詩張之\end{fw}

汀州書法出諸城波磔分明直勁成所好容敎歐九聚後來分取邵庵名\jz{虞伯生自號邵庵因其寫康節詩於一室中也今以相擬}護持神物看連璧篤嗜鄕賢守在楹奮筆我能求法乳石經今見漢熹平

\begin{fw}題湯君東父印譜\end{fw}

君昔年少居浙中摹印故與黃\jz{小松}蔣\jz{山堂}同偶然旁遁近陳\jz{曼生}趙\jz{次閑}拙致獨遜丁龍泓印泥璀錯文字紅置身疑在嘉道中不然何從見此手使石化雪刀如風近來刻者炫秦漢文難急就凡將通君遂歛袖嬾再刻不扶自直麻與蓬君毋小技憎雕蟲自苦後必知良工削觚一集世多重我感所友吳缶翁

\begin{fw}感事\end{fw}

入關但近舞裙塵父死忘讐若楚秦執挺降王甘作長裼裘公子果何人先除異己終無補自壞家居久有因太息武公言不用十年往事計金緡

\begin{fw}爲陸丹林題紅樹室圖\end{fw}

詞人昔有萬紅樹何似霜紅自築堪並世江南黃葉在華涇樓榜說劉三\jz{社友劉三於故里華涇署所居曰黃葉樓可相擬也}

\begin{fw}淞南弔夢圖爲陸丹林作\end{fw}

感舊今爲弔夢人莫敎影事等前塵海天兜率神仙地但恐雙成作後身

\begin{fw}二十九日感述\end{fw}

六日已過七日來百千萬刼閱奔雷憂時每說家多燼無地何言骨已灰兵法我知哀者勝戰場世有幾人回先生強說樓居樂欲舉囊錢換野梅\jz{時思覓盆梅置左右而不得}

\begin{fw}壬申二月十一日\end{fw}

貯淚惟枯目忘憂祗此心死生大事業病候象晴陰生意絪縕得酣眠取次尋求醫問我友積感更彌襟

\begin{fw}壬申三月十六日\end{fw}

兒時食物爭羅致病榻今年事可夸獨惜清明寒食後晨興猶未注新茶

掘參往業成王略我啖其根活病軀今向關東論豪傑尙多張耳與陳餘

\chapter*{病起樓詩}
\begin{fw}病中雜書\end{fw}

積凍知敎雨釀成春初意尙在寒城乞醫却使傳新語飽聽宵來折竹聲

體中不適果何如耐可文成倩子書下筆猶能若風雨汝曹遲速似翁無

當年槃辟笑湯翁\jz{蟄先}托疾佯言體有風不信孱軀眞視此杖行猶賴夾持中

十一時行九百里奔輪尙怨到家遲下車墮軌幾身殉利害相懸在一絲

斷巷新看廣陌通高樓鐙火鬧春紅到家視作尋常日太息家人廢舊風

游舊先期待我來遲敎宵半我方囘朱\jz{炎午}王\jz{饒生}隔巷康\jz{特璋}隣宅爲道生還亦壯哉

瞠視家人悔此行下車幸見壻來迎兒曹一例分佳橘仍笑兒曹有怨聲

鐙照隣樓語夜闌誰哀蹩躃累心肝平生有病如無病今日安心亦大難

大女携孫走柳州大兒廢學動經秋去年往事塡胷在今夜元宵感白頭

何物肝脾竟病人先生久欲斷貪嗔一樓下視鄰家瓦坐閱江南一月春

不論雷雨還風日納坐天光據榻看終笑南中太姸煖晉陽甲起尙春寒

鄰巷爭傳飲彈人故敎魂淚感前塵旁觀任說陳其美我意尤哀宋敎仁

彊起曾能上下樓誤令鍼折血紛流裹創壹意供高臥床簀翻敎七日留

卌四年前舊檢書忽容遮眼共徐徐餖飣强及高齊事一笑眞同獺祭魚

故人日出復時過見說兵塵已奮戈病榭何堪聞戰事江南金幣已無多

去住原如不繫舟獨憐婦稚苦羈留無端忽憶林翁\jz{亮生}語飛絮飛花入石頭

肝膽槎枒久不平草間求活尙治生篋書蕩燼如秦火快意當年擁百城

諛壽文猶可換錢家人失笑在當前不敎更入劉叉手低首昌黎我愧賢

春來兩度卽聞雷夢醒曾疑雨亦來究竟天公成底事萬花取次照人開

靨輔無端益痘痂抱持嬌女令呼爺夢中狂笑何如此醜婦無慙在我家

阿章居校月能歸驎亦書來念所依傭媼有兒時在側先生翻媿問朝饑

日進盤餐喜蔬笋幾疑身已到湖山西塍荒盡園畦地尙辦青鞋日往還\jz{謂紅桕山莊}

寒雪當時憶故人詩成自喜共愁呻舊詩焚盡新篇富終信吾家未盡貧

往感秋冬日在胸不因歲過爲朦朧立春後夜難堅臥致疾原知在睡中

詩人曾賦例無兒四十生男已恨遲今感桐棺闕臨問流傳昔哭邁兒詩

女雛第四歲將周兵定囘杭病不瘳我苦亂離今尙病何憐弱骨葬杭州

江國書來屢問歸不知臥病鐍樓扉阿爺强坐親書答有夢宵來一倂飛

破曉兒嗁我亦醒此時光景本常經道園豪語吾嘗誦身似初孩作柱銘

嗜味時敎負客筵重知張喙自天然流亡不敢言關汴準備長饑已隔年

上車猶恐囘中道行樂還思臥浹旬多謝家人莫忘却此歸眞似未歸人

寶蓋山腰一隧通近聞墮石震羣矇病中經過生餘悸萬事於今半鑿空

朱翁\jz{炎午}許我伴輕裝再作吳游豈可忘獨笑鬚髯同似戟累敎試服婦人方

一病能令一藥除知能病起恃方書何如薄海忘呻楚遠溯方書未著初

無雨風中弔墓門鳳陽往牒遠能言唐朝諸李曾無賴自署玄元幾世孫

欲去耽爲竟日遊得春晴卽視高秋六朝五季吾何憶尤恥江東有仲謀

昌平陵樹亦迷離坏土人間有異詞同是英雄守門戶孫家旁葬紫髯兒

靈谷人爭看牡丹早春先到見苔寒八功德水難重問二硏携歸伴寫官

山西椎魯本無文有客殫勤守舊聞同感革除尋血石獨言朱暈起山紋

臺城未見柳條新記昔來游已隔春笑倩夏郎\jz{明德}留聚影此行恰得四三人

張孔成灰辱井留臨春結綺亦千秋尙矜脚力攀重磴不羡旁人踞上頭

湖上春寒二月同疏梅不見發新叢果然世界能收拾盡納須彌芥子中

試馬春風興好奇病骸亦覺狀支離單衣不敵重裘暖籠袖猶言六月宜

清涼重到正斜陽湖漲能分樹外光掃葉一樓題壁滿來時屢念故人狂

兒如奔犢喜登山侍杖猶能伴我閒歸視柳君\jz{翼謀}求短册朱翁應喜不空還

藥鑪經月不離身幾似襄陽傚半人連日猶能共登陟吾詩下筆亦通神

父書能讀兒稱慧家計雖廉徑未荒太息論斤成捆賣市傭夸說此家藏

四年前檢雲中集頗喜咸同有此才鐙下眼明重見此一時眞覺病懷開

同校唐書每夜深劉書惜未篋中尋此來挾取鴻編去文史非關一代心

說詩久欲逃文字論韻猶期定紐聲十卷花薰分舊刻一時伴我語縱橫

書船歸後贈茲編曾感孫侯\jz{俶仁}重昔賢記取中州文獻在七經樓集後應傳

二汪\jz{辟疆旭初}得遇失王\jz{小湘}黃\jz{季剛}遠陟疲敎眩目光園徑經行問淮水故人重見兩歐陽\jz{竟無法孝}

破曉開門聽鵲鴉朱翁促起共還家一車待發同兒女初日方生萬道霞

閉置車帷本等閒緣江到眼是青山吾能默坐忘羣動京口橫林感二鬟

郭生\jz{詠琴}相送難爲别說與朱翁可破顏去歲江船扶掖我天寒曾送老夫還

並坐何從識笑顰我難平視等無人鏡臺額粉匆匆事應笑衰顏頳作春

南部甘陵我少年汝曹頂髮尙鬈然儘教車吏橫干涉忍笑今爲讓坐筵

塔山已過全無色飽飲眞憐負惠泉却憶雨中聯騎日故人催上蔣家船

停軌無端阻去程蒼然落日向西平投巾鬨語求家給此是江鄕志願兵

下車擁翼賴朱翁掉臂遲行廣衆中門巷到家難自辨羣驚一病到雙瞳

來去春郊趁暖風緋桃時見野籬中此行常潤趨吳郡却少穠青間淺紅

青盲幾作杜子夏一跌猶驚石曼卿記取明曹播佳話朱諸同是兩儒生

姊殯移歸能往祭應憐半體未同枯剡谿一櫂全家去忍淚何堪過故廬

病眸自笑作書遲甥尙追言昨病時母骨得歸衰父健我將慰汝更何詞

鄰牆花樹去年詩遮眼無書靜更宜多謝朱翁猶鄭重勸刊病起一囊詩

縛袴甘爲盾筆書江南猶欲著孱儒與人家國今非幸爲政難言海大魚

我家傳物有空青張目何煩辨醉醒拋却道書忘夙慧平生惜未守黃庭

坐臥惟宜著小樓羣兒來亦厭啁啾一牕晴日曈曨在漸暖先思典敝裘

察見淵魚本不祥近何屬目感昏茫果然甁拂求初地尊者先當禮瞎堂

尋行數墨老無聞絳灌於今本不文慙記少年飛動意自稱大小眼將軍

一雨羣兒出戶難閒庭頓作廣場看笑嗁雜作隨簷溜卻似舟行上下灘

紛然何物更關渠坐厭胡床出厭車除卻金陵留十日竟敎兩月作樓居

壓眼何愁有腦脂賸容匡坐自哦詩寫刊我尙期它日無病休忘病起時

\begin{fw}客言\jz{附錄}\end{fw}

青島本我土他族迭據之屹然東海濱政俗多從夷市政尤大備一一見設施蓬萊閱海市過者忘侏離雜居尙遠控譎哉木屐兒客言十年前道過尤可悲森森列兵衛艦舶非國旗道旁見羣博蹲踞還爭嬉一人突起立屋角遂小遺側目視博塲袴解呈其私旁呼尙强答格磔若有詞急走還其所露穢猶淋漓冒利竟忘恥其國良可知先生聞客言不惜狀以詩此縱舉細事實已破藩籬道溺本有罸法弛何至斯文野國所判禮防愼自持客語雖賓戲我憂無已時同川浴男女裸逐已不奇放誕視爲俗其政將中衰首俯尻益仰倒置皆人爲先生操此義不畏他族嗤

\chapter*{輯佚}
\begin{fw}補遺一首\end{fw}

使我破除殘夜睡讀君别後五年詩縱橫著語成唐律窈窕爲音近楚詞老去江湖當自惜求之流輩已難知樓扉閉雨回燈坐如夢鈞天合樂時

\jz{錄者按此爲諸宗元題黃節蒹葭樓詩之作詩後署癸丑夏五月二十八日諸宗元據黃節詩學詩律講義補入}




\end{document}